\section{Experiments}
\label{exp}

% Auto-generated by scripts/generate_main_tables_by_method.py from tables/main_table.md.
% Requires: booktabs, multirow, graphicx, and xcolor with the table option.
\providecommand{\best}[1]{\cellcolor{yellow!35}\textbf{\textcolor{red}{#1}}}
\providecommand{\second}[1]{\cellcolor{black!10}\underline{#1}}
\begin{table*}[!h]
\centering
\caption{Average results on PEMS03, PEMS04, and TFNSW (part 1 of 3). Each row is one baseline, with MAE, RMSE, and MAPE reported for each dataset. Avg is the mean cumulative-horizon score over prediction lengths from 1 to 12. Chronos-2 uses multivariate graph-grouped inference, and its fine-tuned variant uses per-year low-rank adaptation (LoRA). Values are mean $\pm$ standard deviation when available; deterministic zero-shot runs are shown as $\pm0.00$. MAPE is in percent. Bold and underlining mark the best and second-best non-foundation methods; foundation-model rows are not ranked.}
\label{tab:main_by_method_part1}
\setlength{\tabcolsep}{2.4pt}
\footnotesize
\renewcommand{\arraystretch}{1}
\resizebox{\textwidth}{!}{%
\begin{tabular}{ll|ccc|ccc|ccc}
\toprule
\multirow{2}{*}{\textbf{Category}} & \multirow{2}{*}{\textbf{Baseline}} & \multicolumn{3}{c|}{\textbf{PEMS03}} & \multicolumn{3}{c|}{\textbf{PEMS04}} & \multicolumn{3}{c}{\textbf{TFNSW}} \\
\cmidrule(lr){3-5} \cmidrule(lr){6-8} \cmidrule(lr){9-11}
 & & \textbf{MAE} & \textbf{RMSE} & \textbf{MAPE (\%)} & \textbf{MAE} & \textbf{RMSE} & \textbf{MAPE (\%)} & \textbf{MAE} & \textbf{RMSE} & \textbf{MAPE (\%)} \\
\midrule
\multirow{4}{*}{\textbf{\shortstack[l]{Static\\STGNN\\Backbones}}} & DCRNN & 14.99$_{\pm0.18}$ & 24.90$_{\pm0.19}$ & 28.69$_{\pm0.86}$ & 22.22$_{\pm0.27}$ & 34.58$_{\pm0.32}$ & 20.53$_{\pm0.40}$ & 179.87$_{\pm14.50}$ & 303.16$_{\pm19.20}$ & 211.72$_{\pm42.72}$ \\
 & ASTGNN & 15.15$_{\pm0.10}$ & 25.15$_{\pm0.17}$ & 26.56$_{\pm0.43}$ & 22.45$_{\pm0.09}$ & 34.96$_{\pm0.12}$ & 19.29$_{\pm0.42}$ & 140.15$_{\pm0.82}$ & 247.22$_{\pm1.02}$ & 125.28$_{\pm3.02}$ \\
 & TGCN & 15.88$_{\pm0.10}$ & 30.12$_{\pm0.40}$ & 28.95$_{\pm0.89}$ & 23.60$_{\pm0.37}$ & 38.82$_{\pm0.55}$ & 20.15$_{\pm0.28}$ & 208.33$_{\pm1.48}$ & 390.75$_{\pm3.25}$ & 177.95$_{\pm3.89}$ \\
 & GWN & 15.33$_{\pm0.26}$ & 25.15$_{\pm0.39}$ & 28.01$_{\pm0.71}$ & 23.41$_{\pm0.27}$ & 35.69$_{\pm0.22}$ & 19.94$_{\pm0.55}$ & 170.89$_{\pm3.29}$ & 313.61$_{\pm5.86}$ & 138.49$_{\pm10.57}$ \\
\midrule
\multirow{4}{*}{\textbf{\shortstack[l]{Na\"ive \\Schemes}}} & Pretrain & 31.37$_{\pm1.11}$ & 41.64$_{\pm0.99}$ & 50.04$_{\pm2.80}$ & 153.33$_{\pm6.62}$ & 186.37$_{\pm7.58}$ & 82.25$_{\pm2.99}$ & 262.63$_{\pm10.87}$ & 383.38$_{\pm12.29}$ & 430.83$_{\pm35.95}$ \\
 & Retrain & 14.16$_{\pm0.15}$ & 23.46$_{\pm0.14}$ & 26.29$_{\pm0.58}$ & 21.64$_{\pm0.20}$ & 33.58$_{\pm0.22}$ & 19.15$_{\pm0.30}$ & 138.04$_{\pm0.62}$ & 247.58$_{\pm1.00}$ & 104.00$_{\pm3.59}$ \\
 & Online-NN & 17.47$_{\pm0.16}$ & 28.29$_{\pm0.08}$ & 29.62$_{\pm1.04}$ & 38.46$_{\pm0.27}$ & 51.85$_{\pm0.46}$ & 26.93$_{\pm0.71}$ & 174.36$_{\pm3.31}$ & 318.17$_{\pm5.87}$ & 129.82$_{\pm8.24}$ \\
 & Online-AN & 13.80$_{\pm0.05}$ & 22.89$_{\pm0.11}$ & 24.96$_{\pm0.21}$ & 21.01$_{\pm0.26}$ & 32.73$_{\pm0.31}$ & 17.82$_{\pm0.31}$ & 130.61$_{\pm0.96}$ & 235.35$_{\pm1.95}$ & 99.93$_{\pm3.54}$ \\
\midrule
\multirow{4}{*}{\textbf{\shortstack[l]{Evolving\\Graph}}} & TrafficStream & 14.32$_{\pm0.20}$ & 23.60$_{\pm0.28}$ & 27.63$_{\pm1.67}$ & 21.43$_{\pm0.53}$ & 33.35$_{\pm0.60}$ & 18.83$_{\pm0.59}$ & 145.42$_{\pm2.41}$ & 259.40$_{\pm4.67}$ & 106.92$_{\pm8.01}$ \\
 & PECPM & 14.90$_{\pm0.81}$ & 24.32$_{\pm0.98}$ & 27.78$_{\pm2.30}$ & 25.16$_{\pm3.18}$ & 39.17$_{\pm4.12}$ & 22.29$_{\pm3.69}$ & 174.19$_{\pm22.07}$ & 318.61$_{\pm45.67}$ & 106.24$_{\pm4.44}$ \\
 & STKEC & 14.47$_{\pm0.22}$ & 23.89$_{\pm0.34}$ & 27.58$_{\pm0.64}$ & 22.78$_{\pm1.14}$ & 34.78$_{\pm1.18}$ & 21.54$_{\pm1.53}$ & 147.86$_{\pm1.05}$ & 259.82$_{\pm2.21}$ & 148.01$_{\pm7.69}$ \\
 & EAC & 15.91$_{\pm0.48}$ & 26.40$_{\pm0.88}$ & 30.12$_{\pm0.81}$ & 34.02$_{\pm9.03}$ & 48.20$_{\pm9.87}$ & 39.19$_{\pm17.54}$ & 137.90$_{\pm1.42}$ & 245.35$_{\pm1.93}$ & 119.37$_{\pm14.85}$ \\
\midrule
\multirow{2}{*}{\textbf{\shortstack[l]{Retrieval\\TTC}}} & STRAP & 14.15$_{\pm0.13}$ & 23.37$_{\pm0.22}$ & 26.12$_{\pm0.31}$ & 21.93$_{\pm0.67}$ & 33.78$_{\pm0.64}$ & 19.49$_{\pm0.57}$ & 135.57$_{\pm1.80}$ & 243.47$_{\pm4.15}$ & 104.21$_{\pm3.09}$ \\
 & ST-TTC & 13.98$_{\pm0.14}$ & 23.18$_{\pm0.13}$ & 26.15$_{\pm0.57}$ & 21.33$_{\pm0.20}$ & 33.21$_{\pm0.22}$ & 19.04$_{\pm0.29}$ & 144.41$_{\pm0.72}$ & 257.96$_{\pm1.08}$ & 106.33$_{\pm3.01}$ \\
\midrule
\multirow{5}{*}{\textbf{\shortstack[l]{Static\\Forecasting\\Backbones}}} & STID & 15.14$_{\pm0.14}$ & 27.10$_{\pm0.38}$ & 26.44$_{\pm0.62}$ & 22.43$_{\pm0.18}$ & 35.63$_{\pm0.31}$ & 18.51$_{\pm0.26}$ & 153.63$_{\pm6.20}$ & 307.19$_{\pm10.11}$ & 112.25$_{\pm11.68}$ \\
 & STNorm & 18.20$_{\pm0.47}$ & 28.58$_{\pm0.65}$ & 52.94$_{\pm2.57}$ & 23.99$_{\pm0.62}$ & 36.65$_{\pm0.80}$ & 27.98$_{\pm1.51}$ & 148.78$_{\pm2.33}$ & 275.25$_{\pm3.16}$ & 182.00$_{\pm6.10}$ \\
 & iTrans & 14.21$_{\pm0.13}$ & 23.42$_{\pm0.14}$ & 26.79$_{\pm0.71}$ & 21.22$_{\pm0.11}$ & 33.04$_{\pm0.17}$ & 19.36$_{\pm0.25}$ & 144.04$_{\pm0.82}$ & 257.27$_{\pm2.16}$ & 113.97$_{\pm6.78}$ \\
 & DLinear & 15.73$_{\pm0.06}$ & 26.23$_{\pm0.04}$ & 28.55$_{\pm0.98}$ & 23.58$_{\pm0.18}$ & 36.54$_{\pm0.22}$ & 19.60$_{\pm0.33}$ & 182.20$_{\pm1.03}$ & 318.13$_{\pm1.10}$ & 187.17$_{\pm5.55}$ \\
 & STAEFormer & 12.26$_{\pm0.11}$ & \second{20.58$_{\pm0.12}$} & \second{23.06$_{\pm0.64}$} & 19.38$_{\pm0.94}$ & 30.64$_{\pm1.20}$ & 16.97$_{\pm0.95}$ & 96.12$_{\pm1.67}$ & 190.33$_{\pm2.83}$ & 66.43$_{\pm2.35}$ \\
\midrule
\multirow{6}{*}{\textbf{\shortstack[l]{Times-series\\Forecasting\\Foundation\\Model}}} & Chronos-2 (ZS) & 15.27$_{\pm0.00}$ & 26.98$_{\pm0.00}$ & 25.83$_{\pm0.00}$ & 21.58$_{\pm0.00}$ & 35.22$_{\pm0.00}$ & 17.62$_{\pm0.00}$ & 227.74$_{\pm0.00}$ & 440.18$_{\pm0.00}$ & 120.67$_{\pm0.00}$ \\
 & Chronos-2 (FT) & 13.31$_{\pm0.01}$ & 23.66$_{\pm0.03}$ & 22.80$_{\pm0.03}$ & 18.94$_{\pm0.04}$ & 31.89$_{\pm0.12}$ & 15.25$_{\pm0.01}$ & 160.20$_{\pm0.68}$ & 323.47$_{\pm2.06}$ & 78.26$_{\pm0.28}$ \\
 & TimesFM2.5 (ZS) & 15.91$_{\pm0.00}$ & 27.86$_{\pm0.00}$ & 26.71$_{\pm0.00}$ & 23.49$_{\pm0.00}$ & 37.69$_{\pm0.00}$ & 19.32$_{\pm0.00}$ & 266.77$_{\pm0.00}$ & 487.06$_{\pm0.00}$ & 160.45$_{\pm0.00}$ \\
 & TimesFM2.5 (FT) & 14.16$_{\pm0.01}$ & 25.36$_{\pm0.04}$ & 23.79$_{\pm0.01}$ & 20.92$_{\pm0.02}$ & 35.20$_{\pm0.02}$ & 16.65$_{\pm0.02}$ & 215.22$_{\pm0.33}$ & 420.65$_{\pm0.52}$ & 124.25$_{\pm0.83}$ \\
 & Moirai-2.0 (ZS) & 16.19$_{\pm0.00}$ & 28.00$_{\pm0.00}$ & 26.15$_{\pm0.00}$ & 24.52$_{\pm0.00}$ & 38.87$_{\pm0.00}$ & 19.11$_{\pm0.00}$ & 294.05$_{\pm0.00}$ & 545.09$_{\pm0.00}$ & 137.97$_{\pm0.00}$ \\
 & Moirai-2.0 (FT) & 13.95$_{\pm0.01}$ & 24.51$_{\pm0.03}$ & 24.18$_{\pm0.01}$ & 20.68$_{\pm0.03}$ & 34.42$_{\pm0.09}$ & 16.78$_{\pm0.02}$ & 142.36$_{\pm0.28}$ & 303.78$_{\pm0.69}$ & 62.69$_{\pm0.22}$ \\
\midrule
\multirow{2}{*}{\textbf{Ours}} & A2TTA(STAE) & \best{11.10$_{\pm0.04}$} & \best{18.87$_{\pm0.07}$} & \best{21.46$_{\pm0.18}$} & \best{16.70$_{\pm0.11}$} & \best{27.46$_{\pm0.15}$} & \best{14.89$_{\pm0.05}$} & \best{67.89$_{\pm0.32}$} & \best{147.15$_{\pm0.64}$} & \best{54.66$_{\pm0.33}$} \\
 & A2TTA(OLAN) & \second{12.19$_{\pm0.04}$} & 20.72$_{\pm0.07}$ & 23.43$_{\pm0.27}$ & \second{18.15$_{\pm0.05}$} & \second{29.59$_{\pm0.07}$} & \second{16.41$_{\pm0.28}$} & \second{86.78$_{\pm0.24}$} & \second{176.97$_{\pm0.84}$} & \second{66.01$_{\pm1.49}$} \\
\bottomrule
\end{tabular}%
}
\end{table*}



\subsection{Experimental Settings}

PeMS experiments use the first 31 days of each selected year;
Table~\ref{tab:evo_scale} summarizes the full data archive.
Each year uses a chronological 60\%/20\%/20\% train/validation/test split.
PEMS windows contain 12 five-minute steps and TFNSW windows contain 12 hourly
steps. All methods share the graph snapshots, splits, test windows, and metric
code. Normalization uses training statistics only; residual non-finite inputs
are set to zero after normalization, and invalid targets are masked.

Online methods process one chronological window at a time, match nodes by
metadata sensor IDs, and release each target after 12 steps. A$^2$TTA keeps
512 released windows and updates its global state every 64 windows. Static
methods receive the same windows without test labels. We report cumulative
mean absolute error (MAE), root mean square error (RMSE), and mean absolute
percentage error (MAPE) at horizons 3, 6, and 12. Avg averages the cumulative
scores from horizons 1 to 12, with Avg-MAE as the primary metric. Results use
five paired seeds unless noted otherwise.

A$^2$TTA(OLAN) and A$^2$TTA(STAE) load and freeze the exact per-year
Online-AN and STAEFormer checkpoints, respectively. Other than learning rate,
their adaptation settings are shared. Appendix~\ref{app:hyper_details} lists
the learning rates and sensitivity results.



\subsection{Baselines}

We compare nine static forecasters (DCRNN, ASTGNN, TGCN, GWN, STID, STNorm,
iTrans, DLinear, and STAEFormer); Pretrain, Retrain, Online-NN, and Online-AN;
four evolving-graph methods (TrafficStream, PECPM, STKEC, and EAC); STRAP and
ST-TTC; and zero-shot (ZS) or yearly fine-tuned (FT) versions of Chronos-2,
TimesFM~2.5, and Moirai-2.0.

Each matched A$^2$TTA/backbone pair shares checkpoint seeds, yearly cohorts,
and preprocessing. New-sensor results use the metadata set difference
$\mathcal V_y\setminus\mathcal V_{y-1}$ at every graph-growing transition
with the same five paired seeds. This isolates calibration from checkpoint,
cohort, and node-order differences.

\subsection{Main Results}



\paragraph{Overall performance.}
Results for three representative networks appear in
\Cref{tab:main_by_method_part1}.
A$^2$TTA(STAE) has the lowest MAE, RMSE, and MAPE on all three, reducing
Avg-MAE over frozen STAEFormer by 9.4\%, 13.8\%, and 29.4\%. On TFNSW,
A$^2$TTA(OLAN) reduces MAE from 130.61 to 86.78 (33.6\%) and RMSE from
235.35 to 176.97 (24.8\%).

\paragraph{Comparison with time-series foundation models.}
TSFMs are reference rows and are not included in the non-foundation ranking.
A$^2$TTA(STAE) outperforms every zero-shot and fine-tuned TSFM cell in
\Cref{tab:main_by_method_part1}, while updating only its calibrator after
label release. Appendix~\ref{appendix} reports the other seven PEMS districts
and horizons 3, 6, and 12 for all ten networks
(\cref{tab:main_by_method_part2,tab:main_by_method_part3,tab:main_part1,tab:main_part2}).


% \input{table/main1.tex}
% \input{table/main2.tex}



\paragraph{Forecast horizons.}
Across all 12 steps in \Cref{per_horizon}, A$^2$TTA(STAE) has the lowest error
on PEMS03 and TFNSW. Its mean per-step MAE reductions over frozen STAEFormer
are 8.2\% and 27.3\%. Appendix~\ref{appendix} gives results for all sensors and
the newly added subset on all ten datasets
(\cref{per_horizon_all_app,per_horizon_new_app}).
% \begin{table*}[t]
\centering
\caption{Newly-added-sensor forecast errors, part 1/2 (PEMS03, PEMS04, PEMS05, PEMS06, PEMS07; mean$\pm$std over seeds per dataset). Years where the graph did not grow are excluded. Baselines are grouped and ordered to match \Cref{tab:main_full}. \textbf{Bold}: best, \underline{underline}: second best.}
\label{tab:tsas_new_sensors_part1}
\resizebox{\textwidth}{!}{%
\begin{tabular}{ccc|cccc|cccc|cccc|cc|cccc|c}
\toprule
\multirow{2}{*}{\textbf{Dataset}} & \multirow{2}{*}{\textbf{Metric}} & \multirow{2}{*}{\textbf{Horizon}} & \multicolumn{4}{c|}{\textbf{Static STGNN Backbones}} & \multicolumn{4}{c|}{\textbf{Na\"ive Schemes}} & \multicolumn{4}{c|}{\textbf{Evolving-Graph Continual}} & \multicolumn{2}{c|}{\textbf{Retrieval / TTC}} & \multicolumn{4}{c|}{\textbf{Other Static Backbones}} & \multicolumn{1}{c}{\textbf{Ours}} \\
\cmidrule(lr){4-7} \cmidrule(lr){8-11} \cmidrule(lr){12-15} \cmidrule(lr){16-17} \cmidrule(lr){18-21} \cmidrule(lr){22-22}
 &  &  & \textbf{DCRNN} & \textbf{ASTGNN} & \textbf{TGCN} & \textbf{GWN} & \textbf{Pretrain} & \textbf{Retrain} & \textbf{OL-NN} & \textbf{OL-AN} & \textbf{TrafStm} & \textbf{PECPM} & \textbf{STKEC} & \textbf{EAC} & \textbf{STRAP} & \textbf{ST-TTC} & \textbf{STID} & \textbf{STNorm} & \textbf{iTrans} & \textbf{DLinear} & \textbf{A2TTA} \\
\midrule
\multirow{12}{*}{\textbf{PEMS03}} & \multirow{4}{*}{MAE} & 3 & 19.05$_{\pm2.59}$ & 10.09$_{\pm0.14}$ & 18.49$_{\pm0.11}$ & 10.56$_{\pm0.19}$ & 25.21$_{\pm0.97}$ & \underline{9.35$_{\pm0.79}$} & 9.60$_{\pm0.08}$ & 9.40$_{\pm0.05}$ & 9.71$_{\pm0.15}$ & 10.48$_{\pm0.66}$ & 9.89$_{\pm0.21}$ & 10.67$_{\pm0.13}$ & 30.31$_{\pm1.06}$ & 9.72$_{\pm0.08}$ & 57.53$_{\pm5.42}$ & 12.10$_{\pm0.85}$ & 9.75$_{\pm0.07}$ & 10.21$_{\pm0.06}$ & \textbf{9.22$_{\pm0.06}$} \\
 &  & 6 & 19.27$_{\pm2.61}$ & 10.97$_{\pm0.13}$ & 18.86$_{\pm0.10}$ & 11.17$_{\pm0.14}$ & 25.53$_{\pm0.95}$ & \underline{9.98$_{\pm0.85}$} & 10.24$_{\pm0.11}$ & 10.06$_{\pm0.04}$ & 10.35$_{\pm0.17}$ & 11.08$_{\pm0.64}$ & 10.48$_{\pm0.16}$ & 11.61$_{\pm0.33}$ & 29.65$_{\pm0.72}$ & 10.38$_{\pm0.08}$ & 57.58$_{\pm5.49}$ & 12.95$_{\pm1.05}$ & 10.40$_{\pm0.07}$ & 11.12$_{\pm0.04}$ & \textbf{9.91$_{\pm0.07}$} \\
 &  & 12 & 20.27$_{\pm2.56}$ & 12.96$_{\pm0.13}$ & 19.93$_{\pm0.10}$ & 12.69$_{\pm0.20}$ & 26.46$_{\pm0.75}$ & \underline{11.32$_{\pm0.96}$} & 11.53$_{\pm0.19}$ & 11.40$_{\pm0.07}$ & 11.70$_{\pm0.26}$ & 12.38$_{\pm0.62}$ & 11.84$_{\pm0.13}$ & 13.66$_{\pm0.56}$ & 30.28$_{\pm2.71}$ & 11.78$_{\pm0.10}$ & 58.15$_{\pm5.39}$ & 14.94$_{\pm1.55}$ & 11.81$_{\pm0.07}$ & 13.21$_{\pm0.02}$ & \textbf{11.29$_{\pm0.10}$} \\
 &  & Avg & 19.45$_{\pm2.59}$ & 11.17$_{\pm0.13}$ & 19.01$_{\pm0.10}$ & 11.34$_{\pm0.15}$ & 25.66$_{\pm0.88}$ & \underline{10.10$_{\pm0.86}$} & 10.34$_{\pm0.12}$ & 10.17$_{\pm0.05}$ & 10.47$_{\pm0.19}$ & 11.20$_{\pm0.64}$ & 10.62$_{\pm0.17}$ & 11.81$_{\pm0.32}$ & 30.08$_{\pm0.54}$ & 10.50$_{\pm0.08}$ & 57.60$_{\pm5.44}$ & 13.17$_{\pm1.11}$ & 10.53$_{\pm0.07}$ & 11.33$_{\pm0.04}$ & \textbf{10.02$_{\pm0.07}$} \\
\cmidrule(lr){2-22}
 & \multirow{4}{*}{RMSE} & 3 & 28.63$_{\pm2.49}$ & 16.95$_{\pm0.19}$ & 32.98$_{\pm0.34}$ & 17.53$_{\pm0.30}$ & 32.20$_{\pm1.00}$ & \underline{15.62$_{\pm1.34}$} & 16.13$_{\pm0.14}$ & 15.79$_{\pm0.11}$ & 16.13$_{\pm0.14}$ & 17.13$_{\pm0.83}$ & 16.40$_{\pm0.27}$ & 17.46$_{\pm0.23}$ & 46.63$_{\pm1.77}$ & 16.23$_{\pm0.09}$ & 98.72$_{\pm8.19}$ & 19.43$_{\pm1.03}$ & 16.24$_{\pm0.10}$ & 17.22$_{\pm0.04}$ & \textbf{15.43$_{\pm0.10}$} \\
 &  & 6 & 29.11$_{\pm2.49}$ & 18.76$_{\pm0.16}$ & 33.57$_{\pm0.33}$ & 18.76$_{\pm0.22}$ & 33.25$_{\pm0.95}$ & \underline{16.97$_{\pm1.46}$} & 17.46$_{\pm0.18}$ & 17.15$_{\pm0.12}$ & 17.49$_{\pm0.18}$ & 18.45$_{\pm0.79}$ & 17.71$_{\pm0.21}$ & 19.51$_{\pm0.60}$ & 46.23$_{\pm1.49}$ & 17.63$_{\pm0.11}$ & 98.87$_{\pm8.24}$ & 20.87$_{\pm1.24}$ & 17.54$_{\pm0.08}$ & 19.17$_{\pm0.02}$ & \textbf{16.82$_{\pm0.13}$} \\
 &  & 12 & 31.13$_{\pm2.38}$ & 22.67$_{\pm0.16}$ & 35.30$_{\pm0.30}$ & 21.76$_{\pm0.37}$ & 35.80$_{\pm0.74}$ & \underline{19.64$_{\pm1.67}$} & 19.99$_{\pm0.34}$ & 19.74$_{\pm0.17}$ & 20.16$_{\pm0.38}$ & 21.14$_{\pm0.78}$ & 20.49$_{\pm0.22}$ & 23.73$_{\pm1.09}$ & 49.16$_{\pm5.10}$ & 20.42$_{\pm0.13}$ & 99.61$_{\pm8.16}$ & 24.35$_{\pm1.90}$ & 20.25$_{\pm0.06}$ & 23.36$_{\pm0.01}$ & \textbf{19.50$_{\pm0.21}$} \\
 &  & Avg & 29.45$_{\pm2.47}$ & 19.11$_{\pm0.16}$ & 33.80$_{\pm0.33}$ & 19.10$_{\pm0.25}$ & 33.53$_{\pm0.86}$ & \underline{17.16$_{\pm1.47}$} & 17.62$_{\pm0.21}$ & 17.32$_{\pm0.12}$ & 17.68$_{\pm0.21}$ & 18.65$_{\pm0.80}$ & 17.94$_{\pm0.23}$ & 19.89$_{\pm0.58}$ & 47.18$_{\pm1.30}$ & 17.84$_{\pm0.11}$ & 98.90$_{\pm8.21}$ & 21.27$_{\pm1.33}$ & 17.76$_{\pm0.07}$ & 19.54$_{\pm0.02}$ & \textbf{17.00$_{\pm0.14}$} \\
\cmidrule(lr){2-22}
 & \multirow{4}{*}{MAPE (\%)} & 3 & 74.40$_{\pm27.93}$ & \underline{28.76$_{\pm0.71}$} & 36.64$_{\pm1.86}$ & 32.27$_{\pm1.54}$ & 58.98$_{\pm4.12}$ & 30.11$_{\pm1.67}$ & 29.01$_{\pm0.74}$ & 29.40$_{\pm0.81}$ & 30.85$_{\pm1.23}$ & 33.49$_{\pm3.42}$ & 32.05$_{\pm0.95}$ & 34.80$_{\pm2.15}$ & 76.12$_{\pm17.29}$ & 31.05$_{\pm1.03}$ & 61.26$_{\pm2.88}$ & 55.66$_{\pm9.16}$ & 30.16$_{\pm0.41}$ & 29.48$_{\pm0.28}$ & \textbf{28.60$_{\pm0.41}$} \\
 &  & 6 & 74.81$_{\pm27.86}$ & 30.39$_{\pm0.70}$ & 37.04$_{\pm2.18}$ & 33.27$_{\pm1.67}$ & 57.79$_{\pm3.94}$ & 31.11$_{\pm1.65}$ & \underline{29.77$_{\pm0.89}$} & 30.08$_{\pm0.72}$ & 31.36$_{\pm1.46}$ & 34.15$_{\pm3.37}$ & 32.68$_{\pm1.07}$ & 36.55$_{\pm1.15}$ & 75.32$_{\pm18.52}$ & 32.08$_{\pm0.87}$ & 61.98$_{\pm3.04}$ & 57.86$_{\pm9.83}$ & 31.16$_{\pm0.40}$ & 31.56$_{\pm0.36}$ & \textbf{29.50$_{\pm0.38}$} \\
 &  & 12 & 76.89$_{\pm27.61}$ & 34.42$_{\pm0.76}$ & 40.28$_{\pm2.46}$ & 36.05$_{\pm1.89}$ & 57.02$_{\pm2.81}$ & 33.86$_{\pm1.90}$ & \underline{32.06$_{\pm1.90}$} & 32.36$_{\pm0.86}$ & 33.79$_{\pm2.64}$ & 36.24$_{\pm3.45}$ & 35.23$_{\pm1.34}$ & 41.76$_{\pm1.43}$ & 80.92$_{\pm21.82}$ & 34.91$_{\pm0.55}$ & 63.43$_{\pm2.64}$ & 63.56$_{\pm11.18}$ & 33.78$_{\pm0.18}$ & 37.14$_{\pm0.66}$ & \textbf{31.47$_{\pm0.33}$} \\
 &  & Avg & 75.13$_{\pm27.78}$ & 30.85$_{\pm0.73}$ & 37.72$_{\pm2.13}$ & 33.62$_{\pm1.69}$ & 58.10$_{\pm3.76}$ & 31.44$_{\pm1.70}$ & \underline{30.11$_{\pm1.02}$} & 30.41$_{\pm0.79}$ & 31.78$_{\pm1.67}$ & 34.40$_{\pm3.40}$ & 33.10$_{\pm1.10}$ & 37.22$_{\pm1.30}$ & 76.77$_{\pm18.54}$ & 32.42$_{\pm0.81}$ & 62.05$_{\pm2.87}$ & 58.56$_{\pm9.92}$ & 31.48$_{\pm0.33}$ & 32.27$_{\pm0.34}$ & \textbf{29.68$_{\pm0.38}$} \\
\midrule
\multirow{12}{*}{\textbf{PEMS04}} & \multirow{4}{*}{MAE} & 3 & 99.44$_{\pm10.40}$ & 16.47$_{\pm0.26}$ & 111.23$_{\pm1.42}$ & 18.43$_{\pm0.50}$ & 135.51$_{\pm6.16}$ & 16.12$_{\pm0.18}$ & 15.74$_{\pm0.28}$ & 15.70$_{\pm0.17}$ & \underline{15.66$_{\pm0.24}$} & 19.19$_{\pm4.22}$ & 16.80$_{\pm1.41}$ & 26.85$_{\pm8.49}$ & 91.01$_{\pm22.82}$ & 16.12$_{\pm0.18}$ & 96.85$_{\pm0.92}$ & 18.31$_{\pm0.91}$ & 16.05$_{\pm0.27}$ & 39.79$_{\pm2.85}$ & \textbf{14.75$_{\pm0.15}$} \\
 &  & 6 & 99.58$_{\pm10.36}$ & 18.04$_{\pm0.19}$ & 111.37$_{\pm1.39}$ & 19.25$_{\pm0.45}$ & 135.38$_{\pm5.92}$ & 17.31$_{\pm0.12}$ & \underline{16.81$_{\pm0.36}$} & 16.91$_{\pm0.22}$ & 16.85$_{\pm0.26}$ & 20.30$_{\pm4.08}$ & 17.79$_{\pm1.06}$ & 28.45$_{\pm8.00}$ & 89.69$_{\pm23.61}$ & 17.31$_{\pm0.12}$ & 97.07$_{\pm0.63}$ & 19.49$_{\pm0.77}$ & 17.21$_{\pm0.18}$ & 40.87$_{\pm2.80}$ & \textbf{16.05$_{\pm0.18}$} \\
 &  & 12 & 100.27$_{\pm10.23}$ & 21.56$_{\pm0.16}$ & 111.81$_{\pm1.35}$ & 22.00$_{\pm0.49}$ & 135.52$_{\pm5.96}$ & 19.97$_{\pm0.12}$ & \underline{19.07$_{\pm0.62}$} & 19.40$_{\pm0.26}$ & 19.35$_{\pm0.40}$ & 22.79$_{\pm3.93}$ & 20.43$_{\pm0.96}$ & 31.95$_{\pm6.55}$ & 87.04$_{\pm25.16}$ & 19.97$_{\pm0.12}$ & 97.45$_{\pm0.84}$ & 22.35$_{\pm0.53}$ & 19.82$_{\pm0.14}$ & 43.67$_{\pm2.68}$ & \textbf{18.70$_{\pm0.26}$} \\
 &  & Avg & 99.70$_{\pm10.33}$ & 18.39$_{\pm0.20}$ & 111.44$_{\pm1.39}$ & 19.67$_{\pm0.45}$ & 135.47$_{\pm5.97}$ & 17.58$_{\pm0.14}$ & \underline{17.01$_{\pm0.40}$} & 17.12$_{\pm0.21}$ & 17.07$_{\pm0.28}$ & 20.54$_{\pm4.09}$ & 18.12$_{\pm1.16}$ & 28.80$_{\pm7.75}$ & 89.42$_{\pm23.68}$ & 17.58$_{\pm0.14}$ & 97.22$_{\pm0.78}$ & 19.82$_{\pm0.75}$ & 17.47$_{\pm0.20}$ & 41.20$_{\pm2.81}$ & \textbf{16.27$_{\pm0.19}$} \\
\cmidrule(lr){2-22}
 & \multirow{4}{*}{RMSE} & 3 & 120.58$_{\pm12.39}$ & 25.98$_{\pm0.28}$ & 169.79$_{\pm1.24}$ & 28.09$_{\pm0.48}$ & 168.45$_{\pm7.38}$ & 25.35$_{\pm0.21}$ & \underline{24.84$_{\pm0.19}$} & 24.85$_{\pm0.17}$ & 24.86$_{\pm0.25}$ & 29.72$_{\pm6.04}$ & 26.04$_{\pm1.39}$ & 38.57$_{\pm10.24}$ & 117.27$_{\pm25.98}$ & 25.35$_{\pm0.21}$ & 153.72$_{\pm0.83}$ & 28.28$_{\pm1.15}$ & 25.26$_{\pm0.27}$ & 51.59$_{\pm3.07}$ & \textbf{23.84$_{\pm0.15}$} \\
 &  & 6 & 120.77$_{\pm12.29}$ & 28.76$_{\pm0.21}$ & 169.75$_{\pm1.23}$ & 29.68$_{\pm0.52}$ & 168.62$_{\pm7.06}$ & 27.60$_{\pm0.15}$ & \underline{26.87$_{\pm0.37}$} & 27.05$_{\pm0.23}$ & 27.06$_{\pm0.28}$ & 31.75$_{\pm5.80}$ & 28.04$_{\pm1.01}$ & 41.20$_{\pm9.37}$ & 116.27$_{\pm26.59}$ & 27.60$_{\pm0.15}$ & 153.93$_{\pm0.82}$ & 30.19$_{\pm0.93}$ & 27.40$_{\pm0.17}$ & 54.05$_{\pm3.04}$ & \textbf{26.15$_{\pm0.20}$} \\
 &  & 12 & 121.72$_{\pm12.01}$ & 34.63$_{\pm0.21}$ & 169.73$_{\pm1.22}$ & 34.37$_{\pm0.76}$ & 169.82$_{\pm7.22}$ & 32.22$_{\pm0.20}$ & \underline{30.79$_{\pm0.88}$} & 31.28$_{\pm0.33}$ & 31.37$_{\pm0.51}$ & 36.04$_{\pm5.52}$ & 32.77$_{\pm1.11}$ & 46.82$_{\pm6.92}$ & 114.70$_{\pm28.14}$ & 32.22$_{\pm0.20}$ & 154.07$_{\pm1.02}$ & 34.75$_{\pm0.61}$ & 31.92$_{\pm0.19}$ & 59.49$_{\pm2.96}$ & \textbf{30.59$_{\pm0.35}$} \\
 &  & Avg & 120.93$_{\pm12.24}$ & 29.27$_{\pm0.23}$ & 169.76$_{\pm1.23}$ & 30.31$_{\pm0.52}$ & 168.89$_{\pm7.14}$ & 27.99$_{\pm0.18}$ & \underline{27.14$_{\pm0.44}$} & 27.35$_{\pm0.23}$ & 27.37$_{\pm0.32}$ & 32.12$_{\pm5.82}$ & 28.55$_{\pm1.17}$ & 41.75$_{\pm8.97}$ & 116.16$_{\pm26.70}$ & 27.99$_{\pm0.18}$ & 154.00$_{\pm0.71}$ & 30.69$_{\pm0.90}$ & 27.80$_{\pm0.21}$ & 54.57$_{\pm3.05}$ & \textbf{26.46$_{\pm0.22}$} \\
\cmidrule(lr){2-22}
 & \multirow{4}{*}{MAPE (\%)} & 3 & 494.29$_{\pm62.84}$ & 21.49$_{\pm1.08}$ & 68.08$_{\pm4.13}$ & 23.23$_{\pm1.15}$ & 84.88$_{\pm2.60}$ & 21.86$_{\pm0.80}$ & 20.28$_{\pm1.01}$ & \underline{19.72$_{\pm0.41}$} & 21.33$_{\pm0.70}$ & 25.56$_{\pm5.26}$ & 24.61$_{\pm2.21}$ & 40.84$_{\pm17.97}$ & 70.46$_{\pm17.53}$ & 21.86$_{\pm0.80}$ & 43.91$_{\pm0.78}$ & 33.41$_{\pm3.65}$ & 22.32$_{\pm0.70}$ & 34.47$_{\pm3.29}$ & \textbf{19.35$_{\pm0.34}$} \\
 &  & 6 & 494.94$_{\pm62.44}$ & 23.03$_{\pm1.05}$ & 68.34$_{\pm4.18}$ & 23.92$_{\pm1.21}$ & 84.71$_{\pm2.71}$ & 22.96$_{\pm0.70}$ & 21.20$_{\pm1.27}$ & \underline{20.73$_{\pm0.37}$} & 22.33$_{\pm0.79}$ & 26.80$_{\pm5.30}$ & 25.14$_{\pm1.73}$ & 41.81$_{\pm16.49}$ & 69.02$_{\pm16.08}$ & 22.96$_{\pm0.70}$ & 44.29$_{\pm0.80}$ & 34.86$_{\pm3.53}$ & 23.20$_{\pm0.70}$ & 35.75$_{\pm3.11}$ & \textbf{20.22$_{\pm0.36}$} \\
 &  & 12 & 495.68$_{\pm62.28}$ & 27.12$_{\pm0.90}$ & 69.16$_{\pm4.27}$ & 26.23$_{\pm1.25}$ & 84.46$_{\pm2.82}$ & 25.99$_{\pm0.54}$ & 23.53$_{\pm1.97}$ & \underline{23.51$_{\pm0.42}$} & 25.29$_{\pm1.39}$ & 29.63$_{\pm5.00}$ & 27.74$_{\pm1.50}$ & 45.09$_{\pm15.05}$ & 67.94$_{\pm12.99}$ & 25.99$_{\pm0.54}$ & 45.53$_{\pm0.65}$ & 38.59$_{\pm3.56}$ & 25.55$_{\pm0.82}$ & 39.86$_{\pm2.93}$ & \textbf{22.20$_{\pm0.41}$} \\
 &  & Avg & 494.79$_{\pm62.38}$ & 23.53$_{\pm1.01}$ & 68.45$_{\pm4.21}$ & 24.28$_{\pm1.18}$ & 84.72$_{\pm2.69}$ & 23.35$_{\pm0.67}$ & 21.49$_{\pm1.36}$ & \underline{21.09$_{\pm0.34}$} & 22.75$_{\pm0.90}$ & 27.09$_{\pm5.17}$ & 25.55$_{\pm1.74}$ & 42.17$_{\pm16.61}$ & 69.19$_{\pm15.50}$ & 23.35$_{\pm0.67}$ & 44.51$_{\pm0.65}$ & 35.33$_{\pm3.57}$ & 23.52$_{\pm0.75}$ & 36.38$_{\pm3.11}$ & \textbf{20.43$_{\pm0.37}$} \\
\midrule
\multirow{12}{*}{\textbf{PEMS05}} & \multirow{4}{*}{MAE} & 3 & 8.41$_{\pm0.11}$ & 8.02$_{\pm0.05}$ & 8.24$_{\pm0.13}$ & 8.52$_{\pm0.14}$ & 20.46$_{\pm0.98}$ & 8.02$_{\pm0.09}$ & \underline{7.69$_{\pm0.06}$} & 7.74$_{\pm0.09}$ & 7.87$_{\pm0.18}$ & 8.17$_{\pm0.29}$ & 7.94$_{\pm0.16}$ & 11.48$_{\pm3.43}$ & 23.19$_{\pm5.81}$ & 8.02$_{\pm0.09}$ & 12.61$_{\pm4.42}$ & 9.10$_{\pm0.13}$ & 7.98$_{\pm0.17}$ & 8.13$_{\pm0.05}$ & \textbf{7.50$_{\pm0.08}$} \\
 &  & 6 & 8.64$_{\pm0.08}$ & 8.66$_{\pm0.05}$ & 8.58$_{\pm0.10}$ & 9.02$_{\pm0.15}$ & 20.77$_{\pm0.94}$ & 8.47$_{\pm0.08}$ & \underline{8.07$_{\pm0.06}$} & 8.18$_{\pm0.06}$ & 8.28$_{\pm0.18}$ & 8.62$_{\pm0.27}$ & 8.36$_{\pm0.13}$ & 11.92$_{\pm3.52}$ & 23.27$_{\pm5.85}$ & 8.47$_{\pm0.08}$ & 12.94$_{\pm4.39}$ & 9.46$_{\pm0.14}$ & 8.44$_{\pm0.14}$ & 8.81$_{\pm0.04}$ & \textbf{7.96$_{\pm0.06}$} \\
 &  & 12 & 9.39$_{\pm0.06}$ & 9.97$_{\pm0.07}$ & 9.43$_{\pm0.08}$ & 10.32$_{\pm0.17}$ & 21.62$_{\pm1.00}$ & 9.41$_{\pm0.10}$ & \underline{8.85$_{\pm0.09}$} & 9.04$_{\pm0.05}$ & 9.14$_{\pm0.21}$ & 9.57$_{\pm0.27}$ & 9.25$_{\pm0.13}$ & 12.93$_{\pm3.67}$ & 24.37$_{\pm5.30}$ & 9.41$_{\pm0.10}$ & 13.67$_{\pm4.33}$ & 10.26$_{\pm0.18}$ & 9.38$_{\pm0.12}$ & 10.30$_{\pm0.03}$ & \textbf{8.83$_{\pm0.05}$} \\
 &  & Avg & 8.75$_{\pm0.08}$ & 8.77$_{\pm0.06}$ & 8.68$_{\pm0.10}$ & 9.19$_{\pm0.15}$ & 20.89$_{\pm0.95}$ & 8.55$_{\pm0.09}$ & \underline{8.13$_{\pm0.07}$} & 8.25$_{\pm0.06}$ & 8.36$_{\pm0.19}$ & 8.70$_{\pm0.28}$ & 8.44$_{\pm0.14}$ & 12.02$_{\pm3.52}$ & 23.50$_{\pm5.70}$ & 8.55$_{\pm0.09}$ & 13.02$_{\pm4.40}$ & 9.54$_{\pm0.14}$ & 8.52$_{\pm0.14}$ & 8.95$_{\pm0.04}$ & \textbf{8.02$_{\pm0.06}$} \\
\cmidrule(lr){2-22}
 & \multirow{4}{*}{RMSE} & 3 & 12.50$_{\pm0.16}$ & 12.08$_{\pm0.07}$ & 12.35$_{\pm0.13}$ & 12.58$_{\pm0.13}$ & 25.52$_{\pm1.03}$ & 11.98$_{\pm0.09}$ & \underline{11.52$_{\pm0.10}$} & 11.62$_{\pm0.09}$ & 11.80$_{\pm0.20}$ & 12.17$_{\pm0.43}$ & 11.89$_{\pm0.18}$ & 15.88$_{\pm3.26}$ & 31.12$_{\pm6.96}$ & 11.98$_{\pm0.09}$ & 17.90$_{\pm5.79}$ & 13.15$_{\pm0.21}$ & 11.96$_{\pm0.19}$ & 12.27$_{\pm0.06}$ & \textbf{11.31$_{\pm0.08}$} \\
 &  & 6 & 13.03$_{\pm0.12}$ & 13.25$_{\pm0.08}$ & 13.06$_{\pm0.11}$ & 13.49$_{\pm0.15}$ & 26.09$_{\pm1.03}$ & 12.86$_{\pm0.11}$ & \underline{12.21$_{\pm0.13}$} & 12.48$_{\pm0.06}$ & 12.58$_{\pm0.20}$ & 13.07$_{\pm0.40}$ & 12.70$_{\pm0.16}$ & 16.65$_{\pm3.37}$ & 31.67$_{\pm7.11}$ & 12.86$_{\pm0.11}$ & 18.44$_{\pm5.75}$ & 13.75$_{\pm0.21}$ & 12.87$_{\pm0.17}$ & 13.52$_{\pm0.06}$ & \textbf{12.20$_{\pm0.07}$} \\
 &  & 12 & 14.50$_{\pm0.11}$ & 15.55$_{\pm0.10}$ & 14.68$_{\pm0.13}$ & 15.69$_{\pm0.23}$ & 27.46$_{\pm1.18}$ & 14.65$_{\pm0.19}$ & \textbf{13.51$_{\pm0.19}$} & 14.07$_{\pm0.08}$ & 14.16$_{\pm0.21}$ & 14.87$_{\pm0.40}$ & 14.34$_{\pm0.19}$ & 18.24$_{\pm3.55}$ & 33.65$_{\pm6.61}$ & 14.65$_{\pm0.19}$ & 19.65$_{\pm5.68}$ & 15.02$_{\pm0.27}$ & 14.63$_{\pm0.16}$ & 16.09$_{\pm0.08}$ & \underline{13.80$_{\pm0.09}$} \\
 &  & Avg & 13.22$_{\pm0.13}$ & 13.42$_{\pm0.08}$ & 13.22$_{\pm0.12}$ & 13.75$_{\pm0.16}$ & 26.25$_{\pm1.04}$ & 13.01$_{\pm0.12}$ & \textbf{12.29$_{\pm0.13}$} & 12.58$_{\pm0.06}$ & 12.71$_{\pm0.20}$ & 13.21$_{\pm0.41}$ & 12.84$_{\pm0.17}$ & 16.78$_{\pm3.36}$ & 31.94$_{\pm6.92}$ & 13.01$_{\pm0.12}$ & 18.57$_{\pm5.76}$ & 13.86$_{\pm0.22}$ & 13.00$_{\pm0.17}$ & 13.74$_{\pm0.06}$ & \underline{12.29$_{\pm0.07}$} \\
\cmidrule(lr){2-22}
 & \multirow{4}{*}{MAPE (\%)} & 3 & 33.99$_{\pm1.85}$ & 31.42$_{\pm1.45}$ & 35.07$_{\pm1.23}$ & 31.75$_{\pm1.62}$ & 287.97$_{\pm13.31}$ & 33.01$_{\pm1.29}$ & 30.28$_{\pm1.00}$ & \underline{29.57$_{\pm1.03}$} & 31.84$_{\pm2.12}$ & 34.22$_{\pm3.22}$ & 32.55$_{\pm1.67}$ & 91.33$_{\pm91.36}$ & 72.82$_{\pm24.46}$ & 33.01$_{\pm1.29}$ & 35.51$_{\pm3.14}$ & 52.45$_{\pm6.56}$ & 31.31$_{\pm1.84}$ & 32.31$_{\pm1.71}$ & \textbf{28.86$_{\pm0.78}$} \\
 &  & 6 & 34.37$_{\pm1.78}$ & 33.51$_{\pm1.51}$ & 35.88$_{\pm1.13}$ & 32.97$_{\pm1.44}$ & 288.59$_{\pm12.86}$ & 34.06$_{\pm1.17}$ & 31.34$_{\pm1.21}$ & \underline{30.32$_{\pm0.80}$} & 32.67$_{\pm2.60}$ & 35.22$_{\pm3.06}$ & 33.80$_{\pm1.47}$ & 91.86$_{\pm90.86}$ & 72.97$_{\pm21.57}$ & 34.06$_{\pm1.17}$ & 36.57$_{\pm3.25}$ & 53.80$_{\pm6.50}$ & 32.30$_{\pm1.83}$ & 34.60$_{\pm1.75}$ & \textbf{29.79$_{\pm0.71}$} \\
 &  & 12 & 36.05$_{\pm1.78}$ & 38.05$_{\pm1.43}$ & 38.74$_{\pm1.23}$ & 36.09$_{\pm1.00}$ & 292.23$_{\pm15.31}$ & 36.85$_{\pm1.36}$ & 33.80$_{\pm1.82}$ & \underline{32.43$_{\pm0.38}$} & 34.98$_{\pm3.34}$ & 37.80$_{\pm2.97}$ & 37.21$_{\pm1.53}$ & 96.41$_{\pm89.99}$ & 77.04$_{\pm16.92}$ & 36.85$_{\pm1.36}$ & 39.26$_{\pm3.27}$ & 57.37$_{\pm6.58}$ & 34.83$_{\pm1.61}$ & 40.50$_{\pm1.64}$ & \textbf{31.89$_{\pm0.70}$} \\
 &  & Avg & 34.68$_{\pm1.79}$ & 33.93$_{\pm1.44}$ & 36.32$_{\pm1.14}$ & 33.35$_{\pm1.40}$ & 289.51$_{\pm13.63}$ & 34.40$_{\pm1.24}$ & 31.58$_{\pm1.27}$ & \underline{30.57$_{\pm0.75}$} & 32.95$_{\pm2.62}$ & 35.53$_{\pm3.17}$ & 34.21$_{\pm1.55}$ & 92.72$_{\pm90.68}$ & 73.85$_{\pm20.96}$ & 34.40$_{\pm1.24}$ & 36.91$_{\pm3.18}$ & 54.23$_{\pm6.53}$ & 32.61$_{\pm1.76}$ & 35.31$_{\pm1.70}$ & \textbf{29.99$_{\pm0.72}$} \\
\midrule
\multirow{12}{*}{\textbf{PEMS06}} & \multirow{4}{*}{MAE} & 3 & 14.26$_{\pm2.99}$ & 10.47$_{\pm0.04}$ & 15.63$_{\pm0.10}$ & 11.17$_{\pm0.24}$ & 65.79$_{\pm1.35}$ & 10.13$_{\pm0.03}$ & 9.92$_{\pm0.10}$ & \underline{9.84$_{\pm0.12}$} & 9.98$_{\pm0.09}$ & 12.01$_{\pm3.14}$ & 10.39$_{\pm0.14}$ & 14.55$_{\pm3.45}$ & 45.96$_{\pm8.70}$ & 10.13$_{\pm0.03}$ & 49.20$_{\pm3.83}$ & 12.08$_{\pm0.21}$ & 10.36$_{\pm0.16}$ & 10.76$_{\pm0.07}$ & \textbf{9.58$_{\pm0.12}$} \\
 &  & 6 & 14.52$_{\pm3.01}$ & 11.33$_{\pm0.06}$ & 16.04$_{\pm0.09}$ & 11.87$_{\pm0.20}$ & 66.36$_{\pm1.17}$ & 10.76$_{\pm0.03}$ & 10.52$_{\pm0.13}$ & \underline{10.46$_{\pm0.17}$} & 10.62$_{\pm0.11}$ & 12.55$_{\pm3.06}$ & 10.99$_{\pm0.15}$ & 15.61$_{\pm3.25}$ & 45.52$_{\pm8.25}$ & 10.76$_{\pm0.03}$ & 49.25$_{\pm4.05}$ & 12.66$_{\pm0.18}$ & 10.98$_{\pm0.17}$ & 11.74$_{\pm0.06}$ & \textbf{10.26$_{\pm0.16}$} \\
 &  & 12 & 15.57$_{\pm2.98}$ & 13.10$_{\pm0.13}$ & 17.07$_{\pm0.09}$ & 13.48$_{\pm0.22}$ & 66.88$_{\pm1.23}$ & 12.06$_{\pm0.07}$ & 11.72$_{\pm0.20}$ & \underline{11.69$_{\pm0.24}$} & 11.87$_{\pm0.14}$ & 13.70$_{\pm2.90}$ & 12.30$_{\pm0.17}$ & 17.88$_{\pm2.88}$ & 43.40$_{\pm7.61}$ & 12.06$_{\pm0.07}$ & 49.72$_{\pm3.91}$ & 13.91$_{\pm0.19}$ & 12.32$_{\pm0.21}$ & 13.86$_{\pm0.04}$ & \textbf{11.52$_{\pm0.22}$} \\
 &  & Avg & 14.70$_{\pm3.00}$ & 11.47$_{\pm0.07}$ & 16.16$_{\pm0.09}$ & 12.04$_{\pm0.21}$ & 66.26$_{\pm1.22}$ & 10.86$_{\pm0.04}$ & 10.61$_{\pm0.14}$ & \underline{10.55$_{\pm0.17}$} & 10.71$_{\pm0.11}$ & 12.64$_{\pm3.05}$ & 11.10$_{\pm0.16}$ & 15.82$_{\pm3.21}$ & 45.05$_{\pm8.21}$ & 10.86$_{\pm0.04}$ & 49.34$_{\pm3.94}$ & 12.78$_{\pm0.19}$ & 11.10$_{\pm0.18}$ & 11.93$_{\pm0.06}$ & \textbf{10.34$_{\pm0.16}$} \\
\cmidrule(lr){2-22}
 & \multirow{4}{*}{RMSE} & 3 & 20.17$_{\pm3.25}$ & 16.02$_{\pm0.07}$ & 24.68$_{\pm0.14}$ & 16.88$_{\pm0.41}$ & 81.19$_{\pm1.86}$ & 15.43$_{\pm0.05}$ & 15.16$_{\pm0.16}$ & \underline{14.95$_{\pm0.21}$} & 15.17$_{\pm0.15}$ & 18.87$_{\pm5.96}$ & 15.71$_{\pm0.21}$ & 20.88$_{\pm4.35}$ & 59.79$_{\pm11.67}$ & 15.43$_{\pm0.05}$ & 68.42$_{\pm4.46}$ & 17.65$_{\pm0.26}$ & 15.64$_{\pm0.14}$ & 16.51$_{\pm0.08}$ & \textbf{14.57$_{\pm0.18}$} \\
 &  & 6 & 20.78$_{\pm3.26}$ & 17.58$_{\pm0.11}$ & 25.44$_{\pm0.13}$ & 18.15$_{\pm0.33}$ & 82.22$_{\pm1.68}$ & 16.66$_{\pm0.03}$ & 16.34$_{\pm0.22}$ & \underline{16.16$_{\pm0.27}$} & 16.40$_{\pm0.17}$ & 19.87$_{\pm5.64}$ & 16.92$_{\pm0.21}$ & 22.76$_{\pm3.92}$ & 59.64$_{\pm10.69}$ & 16.66$_{\pm0.03}$ & 68.68$_{\pm4.72}$ & 18.59$_{\pm0.22}$ & 16.80$_{\pm0.16}$ & 18.32$_{\pm0.07}$ & \textbf{15.84$_{\pm0.23}$} \\
 &  & 12 & 22.72$_{\pm3.19}$ & 20.49$_{\pm0.22}$ & 27.18$_{\pm0.13}$ & 20.85$_{\pm0.28}$ & 83.64$_{\pm1.64}$ & 18.95$_{\pm0.10}$ & 18.48$_{\pm0.35}$ & \underline{18.37$_{\pm0.38}$} & 18.65$_{\pm0.23}$ & 21.80$_{\pm5.10}$ & 19.28$_{\pm0.23}$ & 26.57$_{\pm3.15}$ & 57.97$_{\pm9.46}$ & 18.95$_{\pm0.10}$ & 69.49$_{\pm4.56}$ & 20.58$_{\pm0.21}$ & 19.16$_{\pm0.24}$ & 21.83$_{\pm0.06}$ & \textbf{18.12$_{\pm0.33}$} \\
 &  & Avg & 21.06$_{\pm3.24}$ & 17.75$_{\pm0.12}$ & 25.62$_{\pm0.14}$ & 18.39$_{\pm0.34}$ & 82.15$_{\pm1.70}$ & 16.79$_{\pm0.04}$ & 16.45$_{\pm0.23}$ & \underline{16.29$_{\pm0.27}$} & 16.52$_{\pm0.17}$ & 19.99$_{\pm5.62}$ & 17.08$_{\pm0.21}$ & 23.08$_{\pm3.84}$ & 59.15$_{\pm10.67}$ & 16.79$_{\pm0.04}$ & 68.78$_{\pm4.59}$ & 18.77$_{\pm0.23}$ & 16.98$_{\pm0.17}$ & 18.56$_{\pm0.07}$ & \textbf{15.96$_{\pm0.24}$} \\
\cmidrule(lr){2-22}
 & \multirow{4}{*}{MAPE (\%)} & 3 & 58.19$_{\pm29.81}$ & 22.49$_{\pm0.48}$ & 28.46$_{\pm1.67}$ & 24.59$_{\pm1.29}$ & 75.82$_{\pm2.75}$ & 23.24$_{\pm1.45}$ & \underline{21.21$_{\pm0.65}$} & 21.67$_{\pm0.66}$ & 21.90$_{\pm0.49}$ & 27.02$_{\pm7.14}$ & 25.97$_{\pm1.52}$ & 36.18$_{\pm18.89}$ & 67.56$_{\pm5.93}$ & 23.24$_{\pm1.45}$ & 50.16$_{\pm1.68}$ & 33.48$_{\pm2.54}$ & 24.85$_{\pm3.08}$ & 23.70$_{\pm1.11}$ & \textbf{20.80$_{\pm0.64}$} \\
 &  & 6 & 58.57$_{\pm29.86}$ & 23.77$_{\pm0.63}$ & 28.85$_{\pm1.65}$ & 25.64$_{\pm1.36}$ & 75.88$_{\pm3.01}$ & 24.13$_{\pm1.19}$ & \underline{21.92$_{\pm0.57}$} & 22.56$_{\pm0.65}$ & 22.51$_{\pm0.44}$ & 27.78$_{\pm7.06}$ & 26.48$_{\pm1.53}$ & 37.23$_{\pm18.98}$ & 69.33$_{\pm8.74}$ & 24.13$_{\pm1.19}$ & 50.11$_{\pm1.10}$ & 34.58$_{\pm2.73}$ & 25.84$_{\pm3.02}$ & 25.54$_{\pm1.17}$ & \textbf{21.72$_{\pm0.88}$} \\
 &  & 12 & 60.22$_{\pm29.91}$ & 27.09$_{\pm1.04}$ & 30.64$_{\pm1.69}$ & 28.30$_{\pm1.49}$ & 75.47$_{\pm3.01}$ & 26.52$_{\pm0.90}$ & \textbf{23.35$_{\pm0.67}$} & 24.51$_{\pm1.23}$ & 24.22$_{\pm0.51}$ & 29.63$_{\pm6.97}$ & 28.19$_{\pm1.50}$ & 39.35$_{\pm19.29}$ & 71.72$_{\pm17.59}$ & 26.52$_{\pm0.90}$ & 50.99$_{\pm1.07}$ & 37.39$_{\pm3.17}$ & 28.02$_{\pm3.11}$ & 30.02$_{\pm1.43}$ & \underline{23.39$_{\pm1.14}$} \\
 &  & Avg & 58.90$_{\pm29.90}$ & 24.15$_{\pm0.67}$ & 29.16$_{\pm1.65}$ & 25.94$_{\pm1.37}$ & 75.85$_{\pm2.88}$ & 24.41$_{\pm1.19}$ & \underline{22.03$_{\pm0.58}$} & 22.75$_{\pm0.79}$ & 22.72$_{\pm0.43}$ & 27.97$_{\pm7.04}$ & 26.68$_{\pm1.56}$ & 37.39$_{\pm19.08}$ & 69.42$_{\pm9.73}$ & 24.41$_{\pm1.19}$ & 50.35$_{\pm1.36}$ & 34.90$_{\pm2.77}$ & 26.03$_{\pm3.07}$ & 26.05$_{\pm1.22}$ & \textbf{21.82$_{\pm0.85}$} \\
\midrule
\multirow{12}{*}{\textbf{PEMS07}} & \multirow{4}{*}{MAE} & 3 & 63.07$_{\pm5.51}$ & 15.08$_{\pm0.18}$ & 126.15$_{\pm1.70}$ & 16.92$_{\pm1.42}$ & 40.98$_{\pm1.89}$ & 14.53$_{\pm0.20}$ & \textbf{13.57$_{\pm0.10}$} & 13.98$_{\pm0.21}$ & 13.81$_{\pm0.09}$ & 16.10$_{\pm1.56}$ & 13.89$_{\pm0.19}$ & 14.93$_{\pm0.39}$ & 43.42$_{\pm18.18}$ & 14.53$_{\pm0.20}$ & 15.06$_{\pm1.49}$ & 21.62$_{\pm1.24}$ & 14.69$_{\pm0.13}$ & 15.01$_{\pm0.12}$ & \underline{13.68$_{\pm0.13}$} \\
 &  & 6 & 63.36$_{\pm5.45}$ & 16.76$_{\pm0.15}$ & 126.34$_{\pm1.70}$ & 17.81$_{\pm1.18}$ & 41.13$_{\pm1.62}$ & 15.86$_{\pm0.14}$ & \textbf{14.75$_{\pm0.11}$} & 15.39$_{\pm0.22}$ & 15.14$_{\pm0.10}$ & 17.34$_{\pm1.47}$ & 15.25$_{\pm0.17}$ & 15.87$_{\pm0.37}$ & 47.40$_{\pm20.15}$ & 15.86$_{\pm0.14}$ & 15.91$_{\pm1.42}$ & 23.50$_{\pm1.35}$ & 16.03$_{\pm0.14}$ & 16.89$_{\pm0.07}$ & \underline{15.14$_{\pm0.16}$} \\
 &  & 12 & 64.80$_{\pm5.38}$ & 20.43$_{\pm0.16}$ & 126.74$_{\pm1.64}$ & 20.20$_{\pm1.10}$ & 41.26$_{\pm1.68}$ & 18.48$_{\pm0.14}$ & \textbf{16.82$_{\pm0.17}$} & 17.92$_{\pm0.22}$ & \underline{17.54$_{\pm0.19}$} & 19.63$_{\pm1.34}$ & 17.72$_{\pm0.15}$ & 17.77$_{\pm0.44}$ & 58.50$_{\pm27.96}$ & 18.48$_{\pm0.14}$ & 18.38$_{\pm1.52}$ & 27.74$_{\pm1.60}$ & 18.74$_{\pm0.16}$ & 20.97$_{\pm0.06}$ & 17.76$_{\pm0.17}$ \\
 &  & Avg & 63.62$_{\pm5.47}$ & 17.10$_{\pm0.16}$ & 126.38$_{\pm1.68}$ & 18.10$_{\pm1.24}$ & 41.08$_{\pm1.72}$ & 16.04$_{\pm0.15}$ & \textbf{14.84$_{\pm0.12}$} & 15.52$_{\pm0.21}$ & \underline{15.27$_{\pm0.11}$} & 17.47$_{\pm1.47}$ & 15.38$_{\pm0.17}$ & 16.00$_{\pm0.38}$ & 48.98$_{\pm21.82}$ & 16.04$_{\pm0.15}$ & 16.23$_{\pm1.45}$ & 23.95$_{\pm1.37}$ & 16.23$_{\pm0.14}$ & 17.26$_{\pm0.08}$ & 15.28$_{\pm0.15}$ \\
\cmidrule(lr){2-22}
 & \multirow{4}{*}{RMSE} & 3 & 89.39$_{\pm6.03}$ & 24.87$_{\pm0.26}$ & 207.14$_{\pm1.48}$ & 27.20$_{\pm1.90}$ & 51.60$_{\pm1.77}$ & 23.85$_{\pm0.22}$ & \textbf{22.69$_{\pm0.10}$} & 23.21$_{\pm0.23}$ & 22.96$_{\pm0.12}$ & 25.78$_{\pm2.44}$ & 23.10$_{\pm0.21}$ & 24.22$_{\pm0.50}$ & 62.01$_{\pm24.03}$ & 23.85$_{\pm0.22}$ & 25.05$_{\pm2.32}$ & 34.19$_{\pm1.53}$ & 24.09$_{\pm0.08}$ & 25.12$_{\pm0.08}$ & \underline{22.92$_{\pm0.14}$} \\
 &  & 6 & 89.94$_{\pm6.01}$ & 27.95$_{\pm0.22}$ & 207.04$_{\pm1.49}$ & 28.79$_{\pm1.50}$ & 52.70$_{\pm1.41}$ & 26.35$_{\pm0.16}$ & \textbf{24.82$_{\pm0.15}$} & 25.69$_{\pm0.25}$ & \underline{25.39$_{\pm0.15}$} & 28.03$_{\pm2.21}$ & 25.54$_{\pm0.19}$ & 26.02$_{\pm0.46}$ & 69.43$_{\pm27.92}$ & 26.35$_{\pm0.16}$ & 26.46$_{\pm2.25}$ & 37.19$_{\pm1.52}$ & 26.49$_{\pm0.11}$ & 28.68$_{\pm0.07}$ & 25.47$_{\pm0.20}$ \\
 &  & 12 & 92.24$_{\pm5.89}$ & 34.29$_{\pm0.29}$ & 206.82$_{\pm1.46}$ & 32.78$_{\pm1.42}$ & 54.33$_{\pm1.42}$ & 30.92$_{\pm0.21}$ & \textbf{28.33$_{\pm0.25}$} & 29.90$_{\pm0.25}$ & 29.44$_{\pm0.26}$ & 31.97$_{\pm1.92}$ & 29.65$_{\pm0.20}$ & \underline{29.15$_{\pm0.64}$} & 88.34$_{\pm41.51}$ & 30.92$_{\pm0.21}$ & 30.31$_{\pm2.38}$ & 44.26$_{\pm1.57}$ & 31.02$_{\pm0.13}$ & 35.79$_{\pm0.10}$ & 29.79$_{\pm0.20}$ \\
 &  & Avg & 90.32$_{\pm6.01}$ & 28.45$_{\pm0.24}$ & 207.02$_{\pm1.47}$ & 29.22$_{\pm1.59}$ & 52.68$_{\pm1.51}$ & 26.60$_{\pm0.18}$ & \textbf{24.91$_{\pm0.15}$} & 25.84$_{\pm0.24}$ & \underline{25.52$_{\pm0.15}$} & 28.21$_{\pm2.23}$ & 25.68$_{\pm0.20}$ & 26.13$_{\pm0.50}$ & 71.90$_{\pm30.60}$ & 26.60$_{\pm0.18}$ & 26.89$_{\pm2.29}$ & 37.97$_{\pm1.51}$ & 26.76$_{\pm0.10}$ & 29.20$_{\pm0.07}$ & 25.62$_{\pm0.17}$ \\
\cmidrule(lr){2-22}
 & \multirow{4}{*}{MAPE (\%)} & 3 & 227.57$_{\pm91.05}$ & 18.75$_{\pm0.39}$ & 55.74$_{\pm1.64}$ & 20.66$_{\pm1.85}$ & 42.64$_{\pm1.49}$ & 18.95$_{\pm0.87}$ & 17.91$_{\pm0.37}$ & 18.00$_{\pm0.38}$ & \underline{17.45$_{\pm0.64}$} & 20.72$_{\pm1.55}$ & 17.57$_{\pm0.69}$ & 20.65$_{\pm1.03}$ & 55.23$_{\pm18.16}$ & 18.95$_{\pm0.87}$ & 21.64$_{\pm3.06}$ & 39.89$_{\pm2.72}$ & 19.54$_{\pm0.55}$ & 18.08$_{\pm0.57}$ & \textbf{17.14$_{\pm0.35}$} \\
 &  & 6 & 227.77$_{\pm90.51}$ & 20.21$_{\pm0.42}$ & 56.28$_{\pm1.56}$ & 21.45$_{\pm1.68}$ & 41.89$_{\pm1.52}$ & 20.03$_{\pm0.90}$ & 18.78$_{\pm0.43}$ & 18.98$_{\pm0.45}$ & \underline{18.49$_{\pm0.66}$} & 21.65$_{\pm1.56}$ & 18.70$_{\pm0.61}$ & 21.81$_{\pm1.31}$ & 57.38$_{\pm18.64}$ & 20.03$_{\pm0.90}$ & 22.66$_{\pm3.18}$ & 42.65$_{\pm3.25}$ & 20.57$_{\pm0.51}$ & 19.76$_{\pm0.55}$ & \textbf{18.10$_{\pm0.35}$} \\
 &  & 12 & 229.22$_{\pm90.29}$ & 24.00$_{\pm0.40}$ & 57.77$_{\pm1.40}$ & 23.99$_{\pm1.58}$ & 40.94$_{\pm2.06}$ & 22.56$_{\pm1.06}$ & 20.88$_{\pm0.78}$ & 21.24$_{\pm0.55}$ & \underline{20.84$_{\pm0.79}$} & 23.76$_{\pm1.63}$ & 21.32$_{\pm0.74}$ & 24.48$_{\pm1.71}$ & 63.85$_{\pm21.03}$ & 22.56$_{\pm1.06}$ & 26.08$_{\pm3.62}$ & 49.11$_{\pm4.41}$ & 23.04$_{\pm0.47}$ & 24.21$_{\pm0.74}$ & \textbf{20.09$_{\pm0.39}$} \\
 &  & Avg & 228.08$_{\pm90.59}$ & 20.65$_{\pm0.40}$ & 56.47$_{\pm1.53}$ & 21.82$_{\pm1.71}$ & 41.90$_{\pm1.61}$ & 20.30$_{\pm0.91}$ & 19.00$_{\pm0.48}$ & 19.20$_{\pm0.41}$ & \underline{18.73$_{\pm0.69}$} & 21.85$_{\pm1.56}$ & 18.97$_{\pm0.67}$ & 22.09$_{\pm1.19}$ & 58.27$_{\pm19.00}$ & 20.30$_{\pm0.91}$ & 23.17$_{\pm3.24}$ & 43.38$_{\pm3.34}$ & 20.83$_{\pm0.48}$ & 20.31$_{\pm0.60}$ & \textbf{18.26$_{\pm0.36}$} \\
\bottomrule
\end{tabular}%
}
\end{table*}

% \begin{table*}[t]
\centering
\caption{Newly-added-sensor forecast errors, part 2/2 (PEMS08, PEMS10, PEMS11, PEMS12; mean$\pm$std over seeds per dataset). Years where the graph did not grow are excluded. Baselines are grouped and ordered to match \Cref{tab:main_full}. \textbf{Bold}: best, \underline{underline}: second best.}
\label{tab:tsas_new_sensors_part2}
\resizebox{\textwidth}{!}{%
\begin{tabular}{ccc|cccc|cccc|cccc|cc|cccc|c}
\toprule
\multirow{2}{*}{\textbf{Dataset}} & \multirow{2}{*}{\textbf{Metric}} & \multirow{2}{*}{\textbf{Horizon}} & \multicolumn{4}{c|}{\textbf{Static STGNN Backbones}} & \multicolumn{4}{c|}{\textbf{Na\"ive Schemes}} & \multicolumn{4}{c|}{\textbf{Evolving-Graph Continual}} & \multicolumn{2}{c|}{\textbf{Retrieval / TTC}} & \multicolumn{4}{c|}{\textbf{Other Static Backbones}} & \multicolumn{1}{c}{\textbf{Ours}} \\
\cmidrule(lr){4-7} \cmidrule(lr){8-11} \cmidrule(lr){12-15} \cmidrule(lr){16-17} \cmidrule(lr){18-21} \cmidrule(lr){22-22}
 &  &  & \textbf{DCRNN} & \textbf{ASTGNN} & \textbf{TGCN} & \textbf{GWN} & \textbf{Pretrain} & \textbf{Retrain} & \textbf{OL-NN} & \textbf{OL-AN} & \textbf{TrafStm} & \textbf{PECPM} & \textbf{STKEC} & \textbf{EAC} & \textbf{STRAP} & \textbf{ST-TTC} & \textbf{STID} & \textbf{STNorm} & \textbf{iTrans} & \textbf{DLinear} & \textbf{A2TTA} \\
\midrule
\multirow{12}{*}{\textbf{PEMS08}} & \multirow{4}{*}{MAE} & 3 & 63.20$_{\pm14.64}$ & 13.94$_{\pm0.08}$ & 90.82$_{\pm1.78}$ & 15.12$_{\pm0.38}$ & 56.25$_{\pm2.73}$ & 13.52$_{\pm0.08}$ & \underline{12.86$_{\pm0.08}$} & 13.03$_{\pm0.13}$ & 12.97$_{\pm0.12}$ & 15.08$_{\pm0.90}$ & 13.58$_{\pm0.57}$ & 14.37$_{\pm0.25}$ & 38.43$_{\pm9.73}$ & 13.52$_{\pm0.08}$ & 111.42$_{\pm3.40}$ & 17.00$_{\pm0.42}$ & 13.76$_{\pm0.03}$ & 14.06$_{\pm0.04}$ & \textbf{12.72$_{\pm0.09}$} \\
 &  & 6 & 63.50$_{\pm14.64}$ & 15.34$_{\pm0.08}$ & 91.03$_{\pm1.77}$ & 15.98$_{\pm0.25}$ & 56.02$_{\pm2.83}$ & 14.61$_{\pm0.06}$ & \underline{13.94$_{\pm0.11}$} & 14.16$_{\pm0.11}$ & 14.08$_{\pm0.13}$ & 16.04$_{\pm0.85}$ & 14.58$_{\pm0.43}$ & 15.21$_{\pm0.22}$ & 39.83$_{\pm9.50}$ & 14.61$_{\pm0.06}$ & 111.70$_{\pm2.79}$ & 18.02$_{\pm0.34}$ & 14.86$_{\pm0.04}$ & 15.61$_{\pm0.02}$ & \textbf{13.88$_{\pm0.10}$} \\
 &  & 12 & 64.57$_{\pm14.44}$ & 18.35$_{\pm0.12}$ & 91.63$_{\pm1.74}$ & 18.28$_{\pm0.15}$ & 55.80$_{\pm3.30}$ & 16.83$_{\pm0.06}$ & \textbf{15.90$_{\pm0.18}$} & 16.23$_{\pm0.11}$ & 16.12$_{\pm0.17}$ & 17.86$_{\pm0.73}$ & 16.59$_{\pm0.38}$ & 16.86$_{\pm0.21}$ & 43.52$_{\pm9.63}$ & 16.83$_{\pm0.06}$ & 111.84$_{\pm2.57}$ & 20.39$_{\pm0.42}$ & 17.03$_{\pm0.08}$ & 18.95$_{\pm0.04}$ & \underline{15.99$_{\pm0.12}$} \\
 &  & Avg & 63.68$_{\pm14.57}$ & 15.61$_{\pm0.08}$ & 91.11$_{\pm1.77}$ & 16.27$_{\pm0.26}$ & 56.03$_{\pm2.93}$ & 14.79$_{\pm0.07}$ & \underline{14.05$_{\pm0.12}$} & 14.28$_{\pm0.11}$ & 14.20$_{\pm0.14}$ & 16.15$_{\pm0.84}$ & 14.73$_{\pm0.47}$ & 15.31$_{\pm0.22}$ & 40.32$_{\pm9.56}$ & 14.79$_{\pm0.07}$ & 111.83$_{\pm2.82}$ & 18.28$_{\pm0.35}$ & 15.02$_{\pm0.04}$ & 15.91$_{\pm0.02}$ & \textbf{14.00$_{\pm0.10}$} \\
\cmidrule(lr){2-22}
 & \multirow{4}{*}{RMSE} & 3 & 81.62$_{\pm17.23}$ & 22.93$_{\pm0.13}$ & 138.77$_{\pm3.38}$ & 23.91$_{\pm0.47}$ & 65.76$_{\pm3.15}$ & 22.14$_{\pm0.10}$ & \underline{21.30$_{\pm0.09}$} & 21.54$_{\pm0.15}$ & 21.43$_{\pm0.13}$ & 24.01$_{\pm1.20}$ & 22.12$_{\pm0.63}$ & 22.77$_{\pm0.31}$ & 52.13$_{\pm12.73}$ & 22.14$_{\pm0.10}$ & 170.06$_{\pm3.34}$ & 25.96$_{\pm0.52}$ & 22.46$_{\pm0.06}$ & 23.19$_{\pm0.04}$ & \textbf{21.19$_{\pm0.10}$} \\
 &  & 6 & 82.29$_{\pm17.12}$ & 25.57$_{\pm0.13}$ & 139.03$_{\pm3.34}$ & 25.67$_{\pm0.25}$ & 66.21$_{\pm3.21}$ & 24.36$_{\pm0.09}$ & \textbf{23.36$_{\pm0.16}$} & 23.76$_{\pm0.13}$ & 23.64$_{\pm0.16}$ & 25.96$_{\pm1.07}$ & 24.19$_{\pm0.46}$ & 24.47$_{\pm0.28}$ & 55.04$_{\pm12.48}$ & 24.36$_{\pm0.09}$ & 170.32$_{\pm2.96}$ & 27.68$_{\pm0.44}$ & 24.59$_{\pm0.07}$ & 26.15$_{\pm0.03}$ & \underline{23.45$_{\pm0.10}$} \\
 &  & 12 & 84.13$_{\pm16.66}$ & 30.71$_{\pm0.18}$ & 139.63$_{\pm3.25}$ & 29.71$_{\pm0.18}$ & 67.30$_{\pm3.67}$ & 28.47$_{\pm0.12}$ & \textbf{26.88$_{\pm0.32}$} & 27.55$_{\pm0.15}$ & 27.39$_{\pm0.29}$ & 29.45$_{\pm0.85}$ & 27.99$_{\pm0.57}$ & 27.43$_{\pm0.35}$ & 61.94$_{\pm13.05}$ & 28.47$_{\pm0.12}$ & 170.24$_{\pm2.53}$ & 31.57$_{\pm0.64}$ & 28.49$_{\pm0.12}$ & 31.98$_{\pm0.05}$ & \underline{27.31$_{\pm0.13}$} \\
 &  & Avg & 82.53$_{\pm17.02}$ & 25.93$_{\pm0.14}$ & 139.08$_{\pm3.34}$ & 26.08$_{\pm0.28}$ & 66.31$_{\pm3.32}$ & 24.61$_{\pm0.10}$ & \textbf{23.51$_{\pm0.17}$} & 23.92$_{\pm0.14}$ & 23.80$_{\pm0.18}$ & 26.15$_{\pm1.06}$ & 24.41$_{\pm0.55}$ & 24.59$_{\pm0.31}$ & 55.82$_{\pm12.59}$ & 24.61$_{\pm0.10}$ & 170.39$_{\pm2.80}$ & 28.08$_{\pm0.46}$ & 24.82$_{\pm0.07}$ & 26.58$_{\pm0.02}$ & \underline{23.62$_{\pm0.11}$} \\
\cmidrule(lr){2-22}
 & \multirow{4}{*}{MAPE (\%)} & 3 & 316.71$_{\pm89.09}$ & 19.29$_{\pm0.28}$ & 58.18$_{\pm1.88}$ & 21.67$_{\pm1.92}$ & 58.26$_{\pm1.28}$ & 20.17$_{\pm0.82}$ & 18.65$_{\pm0.43}$ & \underline{18.36$_{\pm0.18}$} & 19.24$_{\pm0.78}$ & 22.52$_{\pm1.95}$ & 20.80$_{\pm1.74}$ & 26.09$_{\pm1.75}$ & 56.67$_{\pm13.14}$ & 20.17$_{\pm0.82}$ & 53.03$_{\pm1.94}$ & 36.11$_{\pm4.54}$ & 21.46$_{\pm0.77}$ & 18.99$_{\pm0.56}$ & \textbf{17.31$_{\pm0.36}$} \\
 &  & 6 & 317.48$_{\pm89.04}$ & 20.71$_{\pm0.32}$ & 58.58$_{\pm1.82}$ & 22.51$_{\pm1.91}$ & 57.75$_{\pm1.78}$ & 21.23$_{\pm0.70}$ & 19.77$_{\pm0.47}$ & \underline{19.38$_{\pm0.30}$} & 20.33$_{\pm0.73}$ & 23.54$_{\pm1.91}$ & 21.77$_{\pm1.54}$ & 26.29$_{\pm1.62}$ & 59.45$_{\pm13.27}$ & 21.23$_{\pm0.70}$ & 53.27$_{\pm1.87}$ & 37.52$_{\pm4.50}$ & 22.55$_{\pm0.74}$ & 20.65$_{\pm0.50}$ & \textbf{18.20$_{\pm0.36}$} \\
 &  & 12 & 318.88$_{\pm88.78}$ & 24.18$_{\pm0.54}$ & 60.00$_{\pm1.72}$ & 25.09$_{\pm1.86}$ & 57.48$_{\pm2.77}$ & 23.76$_{\pm0.52}$ & 22.30$_{\pm0.69}$ & \underline{21.86$_{\pm0.39}$} & 22.85$_{\pm0.67}$ & 25.75$_{\pm1.86}$ & 23.98$_{\pm1.05}$ & 27.74$_{\pm2.25}$ & 66.63$_{\pm13.58}$ & 23.76$_{\pm0.52}$ & 54.62$_{\pm1.86}$ & 41.10$_{\pm4.52}$ & 25.07$_{\pm0.73}$ & 24.82$_{\pm0.61}$ & \textbf{20.04$_{\pm0.34}$} \\
 &  & Avg & 317.57$_{\pm88.99}$ & 21.08$_{\pm0.34}$ & 58.82$_{\pm1.84}$ & 22.87$_{\pm1.90}$ & 57.92$_{\pm1.82}$ & 21.50$_{\pm0.66}$ & 20.02$_{\pm0.50}$ & \underline{19.66$_{\pm0.26}$} & 20.59$_{\pm0.70}$ & 23.74$_{\pm1.90}$ & 22.00$_{\pm1.45}$ & 26.54$_{\pm1.50}$ & 60.35$_{\pm13.14}$ & 21.50$_{\pm0.66}$ & 53.59$_{\pm1.75}$ & 37.96$_{\pm4.51}$ & 22.83$_{\pm0.72}$ & 21.14$_{\pm0.52}$ & \textbf{18.36$_{\pm0.35}$} \\
\midrule
\multirow{12}{*}{\textbf{PEMS10}} & \multirow{4}{*}{MAE} & 3 & 10.50$_{\pm0.21}$ & 9.95$_{\pm0.05}$ & 10.98$_{\pm0.07}$ & 10.41$_{\pm0.16}$ & 83.95$_{\pm3.34}$ & 9.69$_{\pm0.06}$ & \textbf{9.39$_{\pm0.09}$} & 9.56$_{\pm0.05}$ & 9.70$_{\pm0.09}$ & 10.45$_{\pm1.00}$ & 9.73$_{\pm0.09}$ & 11.59$_{\pm1.92}$ & 25.66$_{\pm3.68}$ & 9.69$_{\pm0.06}$ & 10.43$_{\pm0.13}$ & 11.42$_{\pm0.26}$ & 9.78$_{\pm0.07}$ & 10.10$_{\pm0.05}$ & \underline{9.41$_{\pm0.04}$} \\
 &  & 6 & 10.74$_{\pm0.17}$ & 10.69$_{\pm0.03}$ & 11.37$_{\pm0.06}$ & 10.85$_{\pm0.17}$ & 84.23$_{\pm3.40}$ & 10.25$_{\pm0.06}$ & \textbf{9.90$_{\pm0.14}$} & 10.13$_{\pm0.07}$ & 10.26$_{\pm0.11}$ & 10.96$_{\pm0.93}$ & 10.26$_{\pm0.09}$ & 12.13$_{\pm2.00}$ & 25.94$_{\pm4.48}$ & 10.25$_{\pm0.06}$ & 10.90$_{\pm0.11}$ & 11.84$_{\pm0.28}$ & 10.32$_{\pm0.07}$ & 10.92$_{\pm0.05}$ & \underline{9.98$_{\pm0.05}$} \\
 &  & 12 & 11.74$_{\pm0.17}$ & 12.23$_{\pm0.02}$ & 12.38$_{\pm0.06}$ & 12.01$_{\pm0.22}$ & 84.43$_{\pm3.39}$ & 11.43$_{\pm0.04}$ & \textbf{10.92$_{\pm0.23}$} & 11.26$_{\pm0.09}$ & 11.44$_{\pm0.15}$ & 12.05$_{\pm0.84}$ & 11.38$_{\pm0.11}$ & 13.42$_{\pm2.24}$ & 29.38$_{\pm9.77}$ & 11.43$_{\pm0.04}$ & 12.12$_{\pm0.13}$ & 12.88$_{\pm0.32}$ & 11.44$_{\pm0.09}$ & 12.71$_{\pm0.05}$ & \underline{11.13$_{\pm0.06}$} \\
 &  & Avg & 10.92$_{\pm0.18}$ & 10.82$_{\pm0.03}$ & 11.50$_{\pm0.06}$ & 10.99$_{\pm0.18}$ & 84.20$_{\pm3.29}$ & 10.36$_{\pm0.05}$ & \textbf{9.98$_{\pm0.15}$} & 10.22$_{\pm0.07}$ & 10.37$_{\pm0.11}$ & 11.06$_{\pm0.93}$ & 10.36$_{\pm0.09}$ & 12.27$_{\pm2.03}$ & 26.73$_{\pm5.48}$ & 10.36$_{\pm0.05}$ & 11.04$_{\pm0.12}$ & 11.96$_{\pm0.27}$ & 10.41$_{\pm0.07}$ & 11.09$_{\pm0.05}$ & \underline{10.07$_{\pm0.05}$} \\
\cmidrule(lr){2-22}
 & \multirow{4}{*}{RMSE} & 3 & 16.84$_{\pm0.44}$ & 16.00$_{\pm0.11}$ & 19.16$_{\pm0.16}$ & 16.54$_{\pm0.19}$ & 86.11$_{\pm3.42}$ & 15.46$_{\pm0.09}$ & \underline{15.03$_{\pm0.15}$} & 15.23$_{\pm0.09}$ & 15.47$_{\pm0.11}$ & 17.13$_{\pm2.26}$ & 15.54$_{\pm0.10}$ & 17.21$_{\pm1.72}$ & 35.95$_{\pm5.42}$ & 15.46$_{\pm0.09}$ & 16.56$_{\pm0.20}$ & 17.70$_{\pm0.40}$ & 15.65$_{\pm0.13}$ & 16.19$_{\pm0.06}$ & \textbf{15.00$_{\pm0.06}$} \\
 &  & 6 & 17.43$_{\pm0.38}$ & 17.37$_{\pm0.08}$ & 19.87$_{\pm0.15}$ & 17.41$_{\pm0.21}$ & 86.77$_{\pm3.55}$ & 16.58$_{\pm0.09}$ & \textbf{16.03$_{\pm0.24}$} & 16.36$_{\pm0.13}$ & 16.57$_{\pm0.16}$ & 18.14$_{\pm2.09}$ & 16.61$_{\pm0.12}$ & 18.26$_{\pm1.75}$ & 37.22$_{\pm7.11}$ & 16.58$_{\pm0.09}$ & 17.44$_{\pm0.16}$ & 18.45$_{\pm0.46}$ & 16.71$_{\pm0.14}$ & 17.71$_{\pm0.06}$ & \underline{16.12$_{\pm0.09}$} \\
 &  & 12 & 19.30$_{\pm0.38}$ & 20.05$_{\pm0.05}$ & 21.60$_{\pm0.14}$ & 19.46$_{\pm0.30}$ & 87.72$_{\pm3.63}$ & 18.82$_{\pm0.06}$ & \textbf{17.92$_{\pm0.43}$} & 18.53$_{\pm0.15}$ & 18.75$_{\pm0.26}$ & 20.14$_{\pm1.80}$ & 18.74$_{\pm0.16}$ & 20.43$_{\pm1.96}$ & 43.59$_{\pm15.75}$ & 18.82$_{\pm0.06}$ & 19.48$_{\pm0.19}$ & 20.22$_{\pm0.56}$ & 18.80$_{\pm0.14}$ & 20.80$_{\pm0.06}$ & \underline{18.35$_{\pm0.12}$} \\
 &  & Avg & 17.71$_{\pm0.40}$ & 17.57$_{\pm0.07}$ & 20.07$_{\pm0.15}$ & 17.63$_{\pm0.22}$ & 86.80$_{\pm3.44}$ & 16.76$_{\pm0.08}$ & \textbf{16.16$_{\pm0.26}$} & 16.51$_{\pm0.12}$ & 16.75$_{\pm0.17}$ & 18.30$_{\pm2.08}$ & 16.78$_{\pm0.12}$ & 18.43$_{\pm1.78}$ & 38.41$_{\pm8.72}$ & 16.76$_{\pm0.08}$ & 17.65$_{\pm0.18}$ & 18.64$_{\pm0.45}$ & 16.87$_{\pm0.14}$ & 17.97$_{\pm0.06}$ & \underline{16.29$_{\pm0.08}$} \\
\cmidrule(lr){2-22}
 & \multirow{4}{*}{MAPE (\%)} & 3 & 32.35$_{\pm0.71}$ & 29.97$_{\pm0.78}$ & 31.64$_{\pm0.79}$ & 31.72$_{\pm1.66}$ & 948.50$_{\pm34.27}$ & 31.31$_{\pm1.00}$ & \textbf{29.29$_{\pm0.39}$} & 30.84$_{\pm0.47}$ & 32.37$_{\pm0.81}$ & 31.02$_{\pm1.72}$ & 31.31$_{\pm0.60}$ & 66.78$_{\pm48.94}$ & 138.67$_{\pm35.18}$ & 31.31$_{\pm1.00}$ & 34.71$_{\pm0.85}$ & 44.44$_{\pm2.05}$ & 30.94$_{\pm1.15}$ & 30.48$_{\pm0.68}$ & \underline{29.40$_{\pm0.32}$} \\
 &  & 6 & 33.04$_{\pm0.67}$ & 31.54$_{\pm0.81}$ & 32.45$_{\pm0.79}$ & 32.65$_{\pm1.62}$ & 946.34$_{\pm33.65}$ & 32.39$_{\pm1.03}$ & \underline{30.34$_{\pm0.53}$} & 31.69$_{\pm0.38}$ & 33.62$_{\pm0.84}$ & 32.01$_{\pm1.68}$ & 32.39$_{\pm0.72}$ & 67.39$_{\pm50.04}$ & 131.51$_{\pm31.92}$ & 32.39$_{\pm1.03}$ & 35.94$_{\pm0.78}$ & 45.66$_{\pm2.03}$ & 31.97$_{\pm1.33}$ & 32.50$_{\pm0.62}$ & \textbf{30.23$_{\pm0.27}$} \\
 &  & 12 & 35.48$_{\pm0.77}$ & 35.40$_{\pm0.87}$ & 35.14$_{\pm0.83}$ & 35.43$_{\pm1.66}$ & 941.79$_{\pm29.53}$ & 35.20$_{\pm1.09}$ & \underline{32.87$_{\pm1.09}$} & 34.31$_{\pm0.49}$ & 37.22$_{\pm0.96}$ & 34.78$_{\pm1.67}$ & 35.01$_{\pm0.81}$ & 71.78$_{\pm53.06}$ & 132.50$_{\pm33.98}$ & 35.20$_{\pm1.09}$ & 39.44$_{\pm0.83}$ & 48.96$_{\pm2.05}$ & 34.53$_{\pm1.67}$ & 37.55$_{\pm0.56}$ & \textbf{32.22$_{\pm0.27}$} \\
 &  & Avg & 33.42$_{\pm0.71}$ & 31.97$_{\pm0.80}$ & 32.85$_{\pm0.79}$ & 33.04$_{\pm1.64}$ & 946.03$_{\pm32.27}$ & 32.75$_{\pm1.02}$ & \underline{30.64$_{\pm0.61}$} & 32.07$_{\pm0.40}$ & 34.11$_{\pm0.80}$ & 32.38$_{\pm1.68}$ & 32.71$_{\pm0.68}$ & 68.33$_{\pm50.45}$ & 134.18$_{\pm32.27}$ & 32.75$_{\pm1.02}$ & 36.39$_{\pm0.79}$ & 46.09$_{\pm2.03}$ & 32.26$_{\pm1.35}$ & 33.09$_{\pm0.60}$ & \textbf{30.45$_{\pm0.27}$} \\
\midrule
\multirow{12}{*}{\textbf{PEMS11}} & \multirow{4}{*}{MAE} & 3 & 52.60$_{\pm9.32}$ & 16.96$_{\pm0.10}$ & 84.02$_{\pm5.17}$ & 18.47$_{\pm0.48}$ & 87.20$_{\pm3.86}$ & 16.41$_{\pm0.21}$ & \underline{15.91$_{\pm0.16}$} & 15.97$_{\pm0.27}$ & 16.27$_{\pm0.19}$ & 24.70$_{\pm3.26}$ & 17.02$_{\pm0.32}$ & 23.34$_{\pm3.67}$ & 56.94$_{\pm10.89}$ & 16.41$_{\pm0.21}$ & 16.42$_{\pm2.92}$ & 20.16$_{\pm0.67}$ & 16.66$_{\pm0.29}$ & 17.22$_{\pm0.16}$ & \textbf{15.26$_{\pm0.17}$} \\
 &  & 6 & 53.24$_{\pm9.23}$ & 18.93$_{\pm0.08}$ & 84.14$_{\pm5.16}$ & 19.65$_{\pm0.48}$ & 88.94$_{\pm3.54}$ & 18.09$_{\pm0.16}$ & 17.52$_{\pm0.21}$ & 17.69$_{\pm0.23}$ & 18.02$_{\pm0.16}$ & 26.11$_{\pm3.10}$ & 18.73$_{\pm0.20}$ & 25.02$_{\pm4.08}$ & 58.28$_{\pm10.35}$ & 18.09$_{\pm0.16}$ & \underline{17.41$_{\pm2.98}$} & 21.69$_{\pm0.82}$ & 18.32$_{\pm0.24}$ & 19.30$_{\pm0.11}$ & \textbf{17.06$_{\pm0.13}$} \\
 &  & 12 & 55.51$_{\pm9.05}$ & 23.17$_{\pm0.15}$ & 84.89$_{\pm5.06}$ & 22.82$_{\pm0.72}$ & 92.59$_{\pm3.19}$ & 21.58$_{\pm0.17}$ & 20.73$_{\pm0.31}$ & 20.99$_{\pm0.20}$ & 21.48$_{\pm0.22}$ & 29.03$_{\pm2.86}$ & 22.20$_{\pm0.22}$ & 29.10$_{\pm4.48}$ & 62.94$_{\pm8.34}$ & 21.58$_{\pm0.17}$ & \textbf{20.16$_{\pm3.22}$} & 24.84$_{\pm1.26}$ & 21.85$_{\pm0.21}$ & 23.94$_{\pm0.11}$ & \underline{20.51$_{\pm0.12}$} \\
 &  & Avg & 53.61$_{\pm9.22}$ & 19.32$_{\pm0.09}$ & 84.29$_{\pm5.14}$ & 20.05$_{\pm0.52}$ & 89.25$_{\pm3.53}$ & 18.39$_{\pm0.17}$ & \underline{17.77$_{\pm0.21}$} & 17.92$_{\pm0.24}$ & 18.29$_{\pm0.19}$ & 26.35$_{\pm3.09}$ & 19.02$_{\pm0.22}$ & 25.48$_{\pm4.02}$ & 59.09$_{\pm9.67}$ & 18.39$_{\pm0.17}$ & 17.77$_{\pm3.02}$ & 21.98$_{\pm0.83}$ & 18.64$_{\pm0.25}$ & 19.76$_{\pm0.13}$ & \textbf{17.31$_{\pm0.15}$} \\
\cmidrule(lr){2-22}
 & \multirow{4}{*}{RMSE} & 3 & 70.68$_{\pm11.75}$ & 26.86$_{\pm0.22}$ & 131.75$_{\pm6.52}$ & 28.77$_{\pm0.66}$ & 110.71$_{\pm3.46}$ & 25.77$_{\pm0.21}$ & \underline{25.09$_{\pm0.15}$} & 25.19$_{\pm0.27}$ & 25.47$_{\pm0.27}$ & 39.42$_{\pm6.00}$ & 26.42$_{\pm0.25}$ & 34.80$_{\pm4.82}$ & 75.77$_{\pm13.30}$ & 25.77$_{\pm0.21}$ & 25.10$_{\pm3.99}$ & 30.89$_{\pm0.73}$ & 26.02$_{\pm0.35}$ & 27.31$_{\pm0.15}$ & \textbf{24.23$_{\pm0.17}$} \\
 &  & 6 & 72.15$_{\pm11.55}$ & 30.73$_{\pm0.19}$ & 131.90$_{\pm6.50}$ & 31.13$_{\pm0.74}$ & 113.75$_{\pm3.17}$ & 29.10$_{\pm0.15}$ & 28.24$_{\pm0.26}$ & 28.48$_{\pm0.25}$ & 28.86$_{\pm0.24}$ & 41.74$_{\pm5.58}$ & 29.83$_{\pm0.15}$ & 37.96$_{\pm5.41}$ & 78.80$_{\pm12.42}$ & 29.10$_{\pm0.15}$ & \textbf{26.91$_{\pm4.18}$} & 33.45$_{\pm0.96}$ & 29.27$_{\pm0.27}$ & 31.45$_{\pm0.09}$ & \underline{27.68$_{\pm0.13}$} \\
 &  & 12 & 76.41$_{\pm11.20}$ & 38.49$_{\pm0.32}$ & 132.76$_{\pm6.38}$ & 36.74$_{\pm1.27}$ & 120.15$_{\pm2.82}$ & 35.53$_{\pm0.25}$ & 34.05$_{\pm0.44}$ & 34.47$_{\pm0.23}$ & 35.11$_{\pm0.38}$ & 46.47$_{\pm4.86}$ & 36.12$_{\pm0.26}$ & 44.86$_{\pm5.89}$ & 87.00$_{\pm10.06}$ & 35.53$_{\pm0.25}$ & \textbf{31.66$_{\pm4.80}$} & 38.62$_{\pm1.73}$ & 35.68$_{\pm0.23}$ & 39.91$_{\pm0.11}$ & \underline{33.95$_{\pm0.17}$} \\
 &  & Avg & 72.73$_{\pm11.53}$ & 31.34$_{\pm0.23}$ & 132.06$_{\pm6.48}$ & 31.71$_{\pm0.83}$ & 114.33$_{\pm3.13}$ & 29.55$_{\pm0.17}$ & 28.60$_{\pm0.25}$ & 28.83$_{\pm0.25}$ & 29.25$_{\pm0.29}$ & 42.11$_{\pm5.54}$ & 30.22$_{\pm0.18}$ & 38.61$_{\pm5.26}$ & 79.96$_{\pm11.58}$ & 29.55$_{\pm0.17}$ & \textbf{27.48$_{\pm4.26}$} & 33.87$_{\pm1.02}$ & 29.75$_{\pm0.29}$ & 32.15$_{\pm0.11}$ & \underline{28.04$_{\pm0.16}$} \\
\cmidrule(lr){2-22}
 & \multirow{4}{*}{MAPE (\%)} & 3 & 248.33$_{\pm60.12}$ & 21.84$_{\pm0.57}$ & 51.33$_{\pm3.72}$ & 23.62$_{\pm1.26}$ & 74.74$_{\pm13.07}$ & 22.70$_{\pm0.44}$ & 21.91$_{\pm0.66}$ & \underline{21.21$_{\pm0.47}$} & 23.02$_{\pm0.43}$ & 32.36$_{\pm6.82}$ & 25.09$_{\pm1.19}$ & 39.73$_{\pm11.18}$ & 72.95$_{\pm11.31}$ & 22.70$_{\pm0.44}$ & 27.97$_{\pm2.98}$ & 36.52$_{\pm2.45}$ & 23.88$_{\pm0.84}$ & 21.85$_{\pm0.34}$ & \textbf{20.48$_{\pm0.32}$} \\
 &  & 6 & 249.14$_{\pm60.20}$ & 23.93$_{\pm0.56}$ & 51.81$_{\pm3.74}$ & 24.86$_{\pm1.32}$ & 77.69$_{\pm14.94}$ & 24.38$_{\pm0.48}$ & 23.34$_{\pm0.70}$ & \underline{22.79$_{\pm0.57}$} & 24.55$_{\pm0.55}$ & 33.87$_{\pm6.63}$ & 26.49$_{\pm0.93}$ & 41.50$_{\pm9.90}$ & 74.59$_{\pm11.17}$ & 24.38$_{\pm0.48}$ & 29.14$_{\pm2.87}$ & 38.85$_{\pm2.76}$ & 25.44$_{\pm0.82}$ & 24.04$_{\pm0.25}$ & \textbf{21.87$_{\pm0.27}$} \\
 &  & 12 & 251.63$_{\pm59.90}$ & 28.83$_{\pm0.58}$ & 53.62$_{\pm3.61}$ & 28.27$_{\pm1.39}$ & 84.80$_{\pm19.50}$ & 28.22$_{\pm0.66}$ & 26.72$_{\pm0.85}$ & \underline{26.37$_{\pm0.65}$} & 28.36$_{\pm1.30}$ & 37.44$_{\pm6.32}$ & 30.18$_{\pm0.72}$ & 46.71$_{\pm9.12}$ & 81.27$_{\pm15.31}$ & 28.22$_{\pm0.66}$ & 33.05$_{\pm2.63}$ & 43.64$_{\pm3.27}$ & 28.98$_{\pm0.87}$ & 29.63$_{\pm0.52}$ & \textbf{24.65$_{\pm0.16}$} \\
 &  & Avg & 249.45$_{\pm60.06}$ & 24.44$_{\pm0.55}$ & 52.09$_{\pm3.69}$ & 25.29$_{\pm1.31}$ & 78.58$_{\pm15.42}$ & 24.78$_{\pm0.50}$ & 23.70$_{\pm0.70}$ & \underline{23.16$_{\pm0.55}$} & 25.02$_{\pm0.69}$ & 34.25$_{\pm6.60}$ & 26.96$_{\pm0.97}$ & 42.20$_{\pm9.95}$ & 75.82$_{\pm11.88}$ & 24.78$_{\pm0.50}$ & 29.72$_{\pm2.79}$ & 39.27$_{\pm2.73}$ & 25.80$_{\pm0.84}$ & 24.72$_{\pm0.32}$ & \textbf{22.09$_{\pm0.26}$} \\
\midrule
\multirow{12}{*}{\textbf{PEMS12}} & \multirow{4}{*}{MAE} & 3 & 30.52$_{\pm9.36}$ & 13.11$_{\pm0.11}$ & 30.59$_{\pm4.06}$ & 13.64$_{\pm0.09}$ & 14.82$_{\pm0.22}$ & 12.54$_{\pm0.13}$ & 12.24$_{\pm0.06}$ & 12.26$_{\pm0.06}$ & 12.24$_{\pm0.08}$ & 13.81$_{\pm1.37}$ & \underline{12.21$_{\pm0.17}$} & 12.88$_{\pm0.23}$ & 32.48$_{\pm7.23}$ & 12.54$_{\pm0.13}$ & 13.43$_{\pm0.37}$ & 15.21$_{\pm0.30}$ & 12.71$_{\pm0.27}$ & 13.12$_{\pm0.08}$ & \textbf{11.66$_{\pm0.13}$} \\
 &  & 6 & 30.89$_{\pm9.30}$ & 14.35$_{\pm0.08}$ & 31.05$_{\pm4.06}$ & 14.47$_{\pm0.11}$ & 15.79$_{\pm0.22}$ & 13.55$_{\pm0.12}$ & \underline{13.14$_{\pm0.07}$} & 13.31$_{\pm0.05}$ & 13.29$_{\pm0.09}$ & 14.74$_{\pm1.28}$ & 13.23$_{\pm0.17}$ & 13.65$_{\pm0.23}$ & 33.02$_{\pm6.56}$ & 13.55$_{\pm0.12}$ & 14.19$_{\pm0.28}$ & 16.20$_{\pm0.34}$ & 13.74$_{\pm0.23}$ & 14.49$_{\pm0.06}$ & \textbf{12.77$_{\pm0.12}$} \\
 &  & 12 & 32.40$_{\pm9.15}$ & 17.13$_{\pm0.06}$ & 32.56$_{\pm4.02}$ & 16.49$_{\pm0.19}$ & 17.68$_{\pm0.20}$ & 15.64$_{\pm0.14}$ & \textbf{14.82$_{\pm0.12}$} & 15.31$_{\pm0.06}$ & 15.30$_{\pm0.16}$ & 16.64$_{\pm1.14}$ & 15.19$_{\pm0.15}$ & 15.29$_{\pm0.37}$ & 35.00$_{\pm5.84}$ & 15.64$_{\pm0.14}$ & 16.29$_{\pm0.25}$ & 18.64$_{\pm0.63}$ & 15.94$_{\pm0.22}$ & 17.49$_{\pm0.10}$ & \underline{14.90$_{\pm0.14}$} \\
 &  & Avg & 31.14$_{\pm9.29}$ & 14.62$_{\pm0.08}$ & 31.28$_{\pm4.05}$ & 14.69$_{\pm0.10}$ & 15.93$_{\pm0.21}$ & 13.72$_{\pm0.13}$ & \underline{13.24$_{\pm0.07}$} & 13.44$_{\pm0.05}$ & 13.42$_{\pm0.09}$ & 14.89$_{\pm1.28}$ & 13.36$_{\pm0.17}$ & 13.79$_{\pm0.25}$ & 33.29$_{\pm6.42}$ & 13.72$_{\pm0.13}$ & 14.45$_{\pm0.30}$ & 16.49$_{\pm0.36}$ & 13.94$_{\pm0.24}$ & 14.77$_{\pm0.07}$ & \textbf{12.92$_{\pm0.13}$} \\
\cmidrule(lr){2-22}
 & \multirow{4}{*}{RMSE} & 3 & 45.00$_{\pm12.05}$ & 22.80$_{\pm0.20}$ & 55.46$_{\pm7.36}$ & 23.14$_{\pm0.27}$ & 23.98$_{\pm0.17}$ & 21.62$_{\pm0.14}$ & 21.31$_{\pm0.10}$ & \underline{21.25$_{\pm0.09}$} & 21.32$_{\pm0.06}$ & 23.70$_{\pm3.00}$ & 21.32$_{\pm0.21}$ & 21.96$_{\pm0.26}$ & 47.90$_{\pm9.87}$ & 21.62$_{\pm0.14}$ & 23.23$_{\pm0.50}$ & 25.07$_{\pm0.38}$ & 21.90$_{\pm0.25}$ & 22.93$_{\pm0.06}$ & \textbf{20.67$_{\pm0.14}$} \\
 &  & 6 & 45.98$_{\pm11.88}$ & 25.43$_{\pm0.15}$ & 56.48$_{\pm7.36}$ & 24.92$_{\pm0.27}$ & 26.05$_{\pm0.19}$ & 23.80$_{\pm0.15}$ & \underline{23.19$_{\pm0.12}$} & 23.43$_{\pm0.08}$ & 23.55$_{\pm0.11}$ & 25.65$_{\pm2.73}$ & 23.47$_{\pm0.19}$ & 23.67$_{\pm0.26}$ & 49.64$_{\pm9.15}$ & 23.80$_{\pm0.15}$ & 24.79$_{\pm0.36}$ & 26.94$_{\pm0.28}$ & 24.04$_{\pm0.21}$ & 25.89$_{\pm0.04}$ & \textbf{22.94$_{\pm0.13}$} \\
 &  & 12 & 48.98$_{\pm11.57}$ & 30.74$_{\pm0.15}$ & 59.28$_{\pm7.23}$ & 28.72$_{\pm0.33}$ & 29.75$_{\pm0.26}$ & 27.83$_{\pm0.22}$ & \textbf{26.41$_{\pm0.21}$} & 27.27$_{\pm0.10}$ & 27.47$_{\pm0.27}$ & 29.31$_{\pm2.35}$ & 27.21$_{\pm0.17}$ & \underline{26.75$_{\pm0.56}$} & 53.94$_{\pm8.91}$ & 27.83$_{\pm0.22}$ & 28.54$_{\pm0.37}$ & 31.31$_{\pm0.71}$ & 28.21$_{\pm0.23}$ & 31.74$_{\pm0.07}$ & 26.94$_{\pm0.18}$ \\
 &  & Avg & 46.38$_{\pm11.85}$ & 25.84$_{\pm0.17}$ & 56.84$_{\pm7.33}$ & 25.25$_{\pm0.26}$ & 26.24$_{\pm0.19}$ & 24.04$_{\pm0.16}$ & \underline{23.32$_{\pm0.13}$} & 23.62$_{\pm0.08}$ & 23.74$_{\pm0.13}$ & 25.87$_{\pm2.73}$ & 23.64$_{\pm0.19}$ & 23.83$_{\pm0.30}$ & 50.01$_{\pm9.12}$ & 24.04$_{\pm0.16}$ & 25.17$_{\pm0.41}$ & 27.41$_{\pm0.30}$ & 24.33$_{\pm0.23}$ & 26.33$_{\pm0.04}$ & \textbf{23.13$_{\pm0.14}$} \\
\cmidrule(lr){2-22}
 & \multirow{4}{*}{MAPE (\%)} & 3 & 189.12$_{\pm81.54}$ & 24.22$_{\pm0.89}$ & 35.89$_{\pm2.38}$ & 24.80$_{\pm0.77}$ & 27.93$_{\pm0.88}$ & 25.12$_{\pm1.29}$ & 23.97$_{\pm0.53}$ & \underline{22.82$_{\pm0.48}$} & 23.83$_{\pm0.63}$ & 27.15$_{\pm2.01}$ & 24.21$_{\pm1.48}$ & 28.87$_{\pm3.67}$ & 98.94$_{\pm55.89}$ & 25.12$_{\pm1.29}$ & 29.91$_{\pm1.34}$ & 49.43$_{\pm6.57}$ & 26.55$_{\pm1.53}$ & 23.91$_{\pm0.76}$ & \textbf{22.15$_{\pm0.24}$} \\
 &  & 6 & 189.83$_{\pm81.85}$ & 26.25$_{\pm0.99}$ & 36.06$_{\pm2.60}$ & 26.07$_{\pm0.94}$ & 28.97$_{\pm1.05}$ & 26.28$_{\pm1.11}$ & 24.98$_{\pm0.53}$ & \underline{23.96$_{\pm0.33}$} & 24.78$_{\pm0.57}$ & 28.11$_{\pm1.60}$ & 25.43$_{\pm1.18}$ & 29.43$_{\pm3.01}$ & 100.19$_{\pm57.31}$ & 26.28$_{\pm1.11}$ & 31.55$_{\pm1.74}$ & 52.52$_{\pm7.86}$ & 27.55$_{\pm1.47}$ & 26.21$_{\pm0.77}$ & \textbf{23.50$_{\pm0.20}$} \\
 &  & 12 & 192.38$_{\pm81.45}$ & 31.60$_{\pm1.38}$ & 39.47$_{\pm2.90}$ & 29.76$_{\pm1.38}$ & 31.53$_{\pm1.19}$ & 29.32$_{\pm0.88}$ & 27.65$_{\pm0.60}$ & \underline{27.15$_{\pm0.39}$} & 27.91$_{\pm0.65}$ & 30.63$_{\pm1.05}$ & 28.60$_{\pm0.78}$ & 31.89$_{\pm3.03}$ & 103.43$_{\pm58.72}$ & 29.32$_{\pm0.88}$ & 36.84$_{\pm2.23}$ & 60.18$_{\pm10.87}$ & 30.38$_{\pm1.29}$ & 32.33$_{\pm1.08}$ & \textbf{26.16$_{\pm0.25}$} \\
 &  & Avg & 190.26$_{\pm81.74}$ & 26.91$_{\pm1.03}$ & 36.87$_{\pm2.61}$ & 26.58$_{\pm1.00}$ & 29.25$_{\pm0.95}$ & 26.64$_{\pm1.12}$ & 25.29$_{\pm0.55}$ & \underline{24.35$_{\pm0.38}$} & 25.20$_{\pm0.58}$ & 28.38$_{\pm1.60}$ & 25.81$_{\pm1.11}$ & 29.83$_{\pm3.17}$ & 100.49$_{\pm57.31}$ & 26.64$_{\pm1.12}$ & 32.31$_{\pm1.73}$ & 53.41$_{\pm8.13}$ & 27.94$_{\pm1.43}$ & 26.98$_{\pm0.82}$ & \textbf{23.70$_{\pm0.21}$} \\
\bottomrule
\end{tabular}%
}
\end{table*}


 \begin{figure}[h!]
    \centering
    \includegraphics[width=0.99\linewidth]{fig/per_horizon_all_pems03_tfnsw_3x2.pdf}
    \caption{All-sensor per-step errors on PEMS03 and TFNSW. Lower is better.}
    \label{per_horizon}
    \Description{Six line plots compare MAE, RMSE, and MAPE over forecast
    steps 1 to 12 on PEMS03 and TFNSW.}
\end{figure}

\subsection{Delayed-Label Sensitivity}

\Cref{fig:delay_sensitivity} uses all ten networks. For PEMS12, Avg-MAE rises by 0.68\% at 524 steps, 0.95\% at 1,036 steps, and 0.65\% without online labels. Across ten networks, the median increases at 524 steps are 0.62\% for A$^2$TTA(OLAN) and 0.63\% for A$^2$TTA(STAE); without online labels, they are 0.44\% and 0.80\%. The sweep changes only the release delay under the one-window protocol.

 \begin{figure}[t]
    \centering
    \includegraphics[width=0.99\linewidth]{fig/delayed_label_common_median_iqr.pdf}
    \caption{Delayed-label sensitivity; bands show the interquartile range.}
    \label{fig:delay_sensitivity}
    \Description{Two curves show the median change in Avg-MAE as labels are
    released later for the Online-AN and STAEFormer variants.}
\end{figure}




\subsection{Ablation Study}

Full A$^2$TTA has the lowest Avg-MAE in seven of eight settings in
\Cref{fig:ablation}. On Online-AN and STAEFormer, respectively, reverting to
the frozen backbone raises error by 6.5\% and 6.7\%, replacing FiLM with a
static affine transform by 3.7\% and 4.7\%, and removing the local clone by
1.2\% and 1.6\%. Freezing FiLM after warm-up is worse in seven settings but
better on STAEFormer with PEMS06. Appendix~\ref{app:ablation_details} defines
the variants and gives the full results.

 \begin{figure}[h!]
    \centering
    \includegraphics[width=0.99\linewidth]{fig/ablation_d_absolute_avg_mae_high_contrast.pdf}
    \caption{Component knockouts on four PEMS datasets and two frozen backbones.}
    \label{fig:ablation}
    \Description{Grouped points compare the full method with four component
    knockouts for Online-AN and STAEFormer on PEMS03 to PEMS06.}
\end{figure}

\subsection{Newly Added Sensors / High-Drift Periods}

 \begin{figure}[h!]
    \centering
    \includegraphics[width=0.99\linewidth]{fig/new_sensor_baselines3.pdf}
    \caption{New-sensor performance.}
    \label{new}
    \Description{Six bar charts compare newly added-sensor MAE, RMSE, and
    MAPE on PEMS07 and TFNSW.}
\end{figure}

For the metadata-defined new-sensor cohort, A$^2$TTA(STAE) is best in five of
six panels in \Cref{new}: all metrics on PEMS07 and MAE and MAPE on TFNSW.
A$^2$TTA(OLAN) has lower TFNSW RMSE (177.20 versus 178.38).
A$^2$TTA(STAE) reduces MAE over frozen STAEFormer by 8.8\% on PEMS07 and
33.7\% on TFNSW. Appendix~\ref{app:new_sensors} gives all ten datasets.







\paragraph{High-drift periods.}

 \begin{figure}[h!]
    \centering
    \includegraphics[width=0.99\linewidth]{fig/drift_gain_main.pdf}
    \caption{Performance across drift-severity quartiles.}
    \label{fig:drift_gain}
    \Description{Points and confidence intervals show MAE reduction across
    quartiles of distribution drift and sensor churn.}
\end{figure}

\Cref{fig:drift_gain} groups paired dataset-year results by drift severity.
A$^2$TTA improves both backbones in every quartile, with all 95\% bootstrap
intervals above zero. Gains peak in the highest sensor-churn quartile at 9.0\%
over Online-AN and 7.9\% over STAEFormer. This legacy-batched diagnostic shows
an association, not a causal effect; Appendix~\ref{app:drift_details} gives
the other metrics and aggregation.




 \begin{figure}[h!]
    \centering
    \includegraphics[width=0.99\linewidth]{fig/case_a2_c_horizontal_PEMS06_2015.pdf}
    \caption{Two-seed case study on PEMS06-2015.}
    \label{fig:case}
     \Description{Two traffic traces compare predictions at short and long
    horizons, and two cohort plots summarize per-sensor MAE changes.}
\end{figure}

\subsection{Case Study}

\Cref{fig:case} examines a high-variance PEMS06-2015 new sensor selected
without using either method's errors. Over seeds 51 and 52, A$^2$TTA tracks
the morning rise more closely, especially at 60 minutes. Across 366 sensors,
it beats Online-AN on 99.7\% and STAEFormer on 98.6\%, reducing mean
per-sensor MAE by 7.29\% and 2.62\%. Appendix~\ref{app:case_details} gives
selection and cohort details.




 \begin{figure}[h!]
    \centering
    \includegraphics[width=0.99\linewidth]{fig/long_term_stability.pdf}
    \caption{Long-term five-seed stability.}
    \label{fig:long_term_stability}
    \Description{Yearly curves show MAE gains over global-only adaptation and
    frozen backbones for six long-running traffic streams.}
\end{figure}

\subsection{Long-Term Adaptation}

The legacy-batched five-seed analysis in \Cref{fig:long_term_stability} covers
six streams over 21 to 25 years. Gains over the frozen backbone are positive
every year; gains over global-only TTA are smaller and less uniform.
Appendix~\ref{app:long_term_details} gives a pre-specified single-seed check
on all ten datasets.




\paragraph{Deployment cost.}
We profile 233 yearly snapshots on an H200. A$^2$TTA(STAE) takes 11.60\,ms
per window versus 10.48\,ms for frozen STAEFormer and raises peak allocated
memory from 10,021 to 10,552\,MiB. It updates 33.2K calibrator parameters on
average versus 2.06M backbone parameters. A$^2$TTA(OLAN) takes 5.35\,ms
versus 4.30\,ms. Appendix~\ref{app:protocol_cost} gives the three-repeat
protocol and additional measurements.
